\documentclass[11pt, a4paper, notitlepage]{report}

\usepackage{amsmath, amssymb, amsthm}
\usepackage{geometry}
\geometry{margin=1in}
\usepackage{hyperref}
\hypersetup{colorlinks=true, linkcolor=blue!60!black, urlcolor=blue!60!black}
\usepackage{enumitem}
\usepackage{algorithm}
\usepackage{algpseudocode}
\usepackage{listings}
\usepackage{xcolor}
\usepackage[most]{tcolorbox}
\tcbuselibrary{theorems}
\usepackage{titlesec}
\usepackage{booktabs}
\usepackage{array}
\setcounter{secnumdepth}{3}

% ── Highlight boxes ───────────────────────────────────────────────────────────
\newtcolorbox{keyresult}[1][]{
  enhanced, colback=blue!5!white, colframe=blue!60!black,
  fonttitle=\bfseries, title={Key Result}, #1
}
\newtcolorbox{examtip}[1][]{
  enhanced, colback=green!5!white, colframe=green!50!black,
  fonttitle=\bfseries, title={Exam Tip}, #1
}
\newtcolorbox{warning}[1][]{
  enhanced, colback=red!5!white, colframe=red!60!black,
  fonttitle=\bfseries, title={Warning}, #1
}

% ── Theorem environments ──────────────────────────────────────────────────────
\newtheorem{theorem}{Theorem}[section]
\newtheorem{lemma}[theorem]{Lemma}
\newtheorem{corollary}[theorem]{Corollary}
\newtheorem{proposition}[theorem]{Proposition}
\theoremstyle{definition}
\newtheorem{definition}[theorem]{Definition}
\newtheorem{example}[theorem]{Example}
\theoremstyle{remark}
\newtheorem{remark}[theorem]{Remark}

% ── Code listings ─────────────────────────────────────────────────────────────
\lstset{
  basicstyle=\ttfamily\small,
  keywordstyle=\color{blue},
  commentstyle=\color{gray},
  breaklines=true,
  frame=single
}

% ── Suppress forced page break in \tableofcontents ────────────────────────────
\makeatletter
\renewcommand\tableofcontents{%
  \section*{\contentsname}%
  \@starttoc{toc}%
}
\makeatother

% ── Chapter title formatting ──────────────────────────────────────────────────
\titleformat{\chapter}[display]
  {\normalfont\LARGE\bfseries}
  {\chaptertitlename\ \thechapter}{12pt}{\Huge}
\titlespacing*{\chapter}{0pt}{20pt}{30pt}

\title{\textbf{MH1301: Discrete Mathematics}\\[0.5em]
       \large Weeks 1--7 Lecture Notes}
\author{NTU --- AY 2025/26}
\date{}

\begin{document}
\maketitle
\tableofcontents
\newpage

% =============================================================================
\chapter{Basics of Counting}
\label{chap:counting}
% =============================================================================

\section{Set Theory Recap}
% =============================================================================

This handout begins with a brief review of set theory from MH1300, which provides the
foundational vocabulary needed for counting.

\subsection{Sets and Notation}

\begin{definition}[Set]
\label{def:set}
A \textbf{set} is an unordered collection of objects, called \emph{elements}. Two sets are
equal if and only if they have exactly the same elements.
\end{definition}

We write $a \in A$ to denote that $a$ is an element of the set $A$, and $a \notin A$ to denote
that $a$ is not an element of $A$. The \textbf{axiom of extension} states that a set is
completely determined by its elements; the order in which elements are listed is irrelevant, and
repeated listings do not change the set. For example,
\[
  \{a, b, c\} = \{b, c, a\} = \{c, a, b\} = \{a, b, c, a\}.
\]

Sets are typically denoted by uppercase letters such as $A$, $B$, $C$. Elements can be listed
explicitly using \textbf{set-roster notation}, e.g.\ $\{1, 2, 3\}$, $\{1, 2, 3, \ldots, 50\}$,
$\{5, 10, 15, 20, \ldots\}$. The symbol ``$\ldots$'' is called an \emph{ellipsis} and is read
``and so forth.''

The standard number sets and their symbols are:

\begin{center}
\begin{tabular}{cl}
\toprule
\textbf{Symbol} & \textbf{Set} \\
\midrule
$\mathbb{N}$ & Set of natural numbers \\
$\mathbb{Z}$ & Set of integers \\
$\mathbb{Q}$ & Set of rational numbers (quotients of integers) \\
$\mathbb{R}$ & Set of real numbers \\
$\mathbb{C}$ & Set of complex numbers \\
\bottomrule
\end{tabular}
\end{center}

In particular, $\mathbb{N} = \{0, 1, 2, 3, \ldots\}$, $\mathbb{Z} = \{\ldots, -2, -1, 0, 1, 2, \ldots\}$,
and $\mathbb{Z}^{\text{nonneg}} = \mathbb{N}$.

A set can also be specified using \textbf{set-builder notation}: given a set $S$ and a property
$P(x)$, the set $\{x \in S \mid P(x)\}$ is read ``the set of all $x$ in $S$ such that $P(x)$
holds.'' For example,
\[
  \mathbb{Z}^+ = \{x \in \mathbb{Z} \mid x > 0\}.
\]

\subsection{Subsets and Cardinality}

\begin{definition}[Subset]
\label{def:subset}
If $A$ and $B$ are sets, then $A$ is called a \textbf{subset} of $B$, written $A \subseteq B$,
if and only if every element of $A$ is also an element of $B$. Symbolically:
\[
  A \subseteq B \iff \forall x,\ (x \in A \Rightarrow x \in B).
\]
We also say ``$A$ is contained in $B$'' or ``$B$ contains $A$.'' The negation $A \not\subseteq B$
means there exists an element $x$ such that $x \in A$ but $x \notin B$.
\end{definition}

\begin{definition}[Proper Subset]
$A$ is a \textbf{proper subset} of $B$, written $A \subsetneq B$, if every element of $A$ is
in $B$ but there is at least one element of $B$ that is not in $A$. The relationship is
analogous to $<$ versus $\leq$ for numbers.
\end{definition}

\begin{definition}[Cardinality]
\label{def:cardinality}
The \textbf{size} or \textbf{cardinality} of a set $S$, denoted $|S|$, is the number of
distinct elements in $S$.
\end{definition}

\begin{example}
What is $|\{a, b, c, a\}|$? Since $\{a, b, c, a\} = \{a, b, c\}$ (repeated elements do not
count twice), the cardinality is $3$.
\end{example}

Counting, as studied in this course, is essentially the task of determining the cardinality of
a given set.

\subsection{Sequences}

\begin{definition}[Sequence]
\label{def:sequence}
A \textbf{sequence} is an \emph{ordered} collection of elements (also called
components or entries), not necessarily distinct.
\end{definition}

Unlike sets, sequences are sensitive to order and repetition. For example, $(a, b, c)$,
$(b, a, c)$, and $(a, b, a)$ are three distinct sequences. The sequences $(a, b, c)$ and
$(b, a, c)$ are \emph{not} equal, even though $\{a, b, c\} = \{b, a, c\}$ as sets.

% =============================================================================
\section{The Product Rule}
% =============================================================================

\subsection{Statement and Intuition}

\begin{keyresult}[title={Key Result: Product Rule}]
Given a set $S$ of sequences of length $k$, where there are $n_1$ possible first entries,
$n_2$ possible second entries for each first entry, $n_3$ possible third entries for each
combination of first and second entries, and so on, the total size of $S$ is
\[
  |S| = n_1 \cdot n_2 \cdot n_3 \cdots n_k.
\]
\end{keyresult}

The product rule applies whenever a process consists of $k$ independent sequential choices.
Each ``box'' (position in the sequence) contributes a multiplicative factor equal to the number
of options available at that stage.

\subsection{Examples}

\begin{example}[Dice and Coin]
A six-sided die and a two-sided coin (H or T) are tossed simultaneously. The set of all
possible outcomes is
\[
  \{(1,H),(2,H),(3,H),(4,H),(5,H),(6,H),(1,T),(2,T),(3,T),(4,T),(5,T),(6,T)\}.
\]
There are 6 choices for the die and 2 for the coin, so by the product rule there are
$6 \times 2 = 12$ outcomes.
\end{example}

\begin{example}[License Plates]
\label{ex:license}
How many different license plates are available if each plate contains a sequence of 3 letters
followed by 3 digits?

\textit{Solution.} There are 26 choices for each of the three letter positions and 10 choices
for each of the three digit positions. By the product rule:
\[
  26 \cdot 26 \cdot 26 \cdot 10 \cdot 10 \cdot 10 = 26^3 \cdot 10^3 = 17{,}576{,}000.
\]
\end{example}

\begin{example}[Permutations of an $n$-element set]
\label{ex:permutation}
A \textbf{permutation} of a set $S$ is a sequence that contains every element of $S$ exactly
once. For example, the permutations of $\{1, 2, 3\}$ are:
$(1,2,3),\ (1,3,2),\ (2,1,3),\ (2,3,1),\ (3,1,2),\ (3,2,1).$

How many permutations does an $n$-element set have? The first entry can be any of the $n$
elements; the second can be any of the remaining $n-1$ elements; the third any of the $n-2$
remaining, and so on. By the product rule:
\[
  n \cdot (n-1) \cdot (n-2) \cdots 1 = n!.
\]
\end{example}

% =============================================================================
\section{Functions and Counting}
% =============================================================================

\subsection{Functions}

\begin{definition}[Function]
\label{def:function}
Let $X$ and $Y$ be sets. A \textbf{function} $f: X \to Y$ is a relation such that for every
element $x \in X$, there is \emph{exactly one} element $y \in Y$ with $f(x) = y$.
\end{definition}

For example, $f(x) = x^2$ defines a function $f: \mathbb{R} \to \mathbb{R}$, since each real
number $x$ maps to a unique value $x^2$. By contrast, $f(x) = \pm\sqrt{x}$ is \emph{not} a
function, because for $x = 4$ it would assign both $+2$ and $-2$.

\begin{definition}[Surjection, Injection, Bijection]
\label{def:bij}
Let $f: X \to Y$.
\begin{itemize}
  \item $f$ is \textbf{surjective} (onto) if every element of $Y$ is mapped from \emph{at least
    one} element of $X$.
  \item $f$ is \textbf{injective} (one-to-one) if every element of $Y$ is mapped from \emph{at
    most one} element of $X$.
  \item $f$ is \textbf{bijective} if it is both surjective and injective, i.e.\ every element
    of $Y$ is mapped from \emph{exactly one} element of $X$.
\end{itemize}
\end{definition}

\subsection{Mapping Rules}

\begin{theorem}[Mapping Rules]
\label{thm:mapping}
Let $X$ and $Y$ be finite sets and $f: X \to Y$ a function.
\begin{itemize}
  \item If $f$ is surjective, then $|X| \geq |Y|$.
  \item If $f$ is injective, then $|X| \leq |Y|$.
  \item If $f$ is bijective, then $|X| = |Y|$.
\end{itemize}
\end{theorem}

\begin{keyresult}[title={Key Result: Bijections and Counting}]
To count $|Y|$, it suffices to find a bijection $f: X \to Y$ and count $|X|$ instead. This is
the fundamental technique of \emph{counting by bijection}: reduce an unfamiliar counting
problem to a familiar one via a one-to-one correspondence.
\end{keyresult}

\begin{warning}[title={Warning: Over- and Under-Counting}]
When constructing a bijection, verify carefully that every element of $Y$ is counted
\emph{exactly once}. Over-counting (surjective but not injective) and under-counting
(injective but not surjective) are the two most common errors.
\end{warning}

\subsection{Counting Functions and Sequences}

\begin{example}[Counting functions]
How many functions are there from a set $X$ with $m$ elements to a set $Y$ with $n$ elements?

\textit{Solution.} Specifying a function $f: X \to Y$ means choosing, for each of the $m$
elements of $X$, one of the $n$ elements of $Y$. By the product rule:
\[
  \underbrace{n \cdot n \cdots n}_{m \text{ times}} = n^m.
\]
\textbf{Note:} The answer is $n^m$, not $m^n$.
\end{example}

\begin{examtip}[title={Exam Tip: $n^m$ vs $m^n$}]
When counting functions from $X$ (size $m$) to $Y$ (size $n$), the answer is $n^m$: for each
element of the \emph{domain} ($m$ elements), there are $n$ choices in the \emph{codomain}.
Do not confuse with $m^n$.
\end{examtip}

\begin{example}[Binary strings]
How many $n$-bit binary sequences are there?

\textit{Solution.} Each bit has 2 possible values (0 or 1). By the product rule, the total
number of $n$-bit sequences is $2^n$. (This equals the number of functions from an $n$-element
set to $\{0, 1\}$.)
\end{example}

\begin{example}[Subsets via bijection]
\label{ex:subsets}
How many different subsets does $S = \{1, 2, \ldots, n\}$ have?

\textit{Solution.} We construct a bijection between subsets of $S$ and $n$-bit binary
sequences. Given a subset $A \subseteq S$, define the corresponding sequence
$(x_1, x_2, \ldots, x_n)$ by:
\[
  x_i = \begin{cases} 1 & \text{if } i \in A, \\ 0 & \text{if } i \notin A. \end{cases}
\]
For example, with $n = 7$, the subset $\{3, 4, 6\}$ maps to $(0, 0, 1, 1, 0, 1, 0)$. Every
subset maps to a unique binary sequence and every binary sequence corresponds to a unique
subset, so this is a bijection. Since there are $2^n$ binary sequences of length $n$, there
are also $2^n$ subsets of $S$.
\end{example}

% =============================================================================
\section{The Sum Rule}
% =============================================================================

\begin{keyresult}[title={Key Result: Sum Rule}]
If $A_1, A_2, \ldots, A_k$ are \textbf{pairwise disjoint} sets, then
\[
  |A_1 \cup A_2 \cup \cdots \cup A_k| = |A_1| + |A_2| + \cdots + |A_k|.
\]
\end{keyresult}

The sum rule applies when the elements to be counted fall into mutually exclusive categories.

\begin{example}[Passwords]
Passwords on a computer are 6 to 8 characters long, where each character is either an
uppercase letter (26 choices) or a digit (10 choices). How many possible passwords are there?

\textit{Solution.} Let $X_6$, $X_7$, $X_8$ denote the sets of valid passwords of length 6,
7, and 8 respectively. These sets are pairwise disjoint, so by the sum rule:
\[
  |X| = |X_6| + |X_7| + |X_8|.
\]
By the product rule (36 choices per character):
\[
  |X_6| = 36^6, \quad |X_7| = 36^7, \quad |X_8| = 36^8.
\]
Therefore the total is $36^6 + 36^7 + 36^8 = 2{,}901{,}650{,}853{,}888 \approx 2.9 \times 10^{12}$.
\end{example}

\begin{warning}[title={Warning: The Sum Rule Requires Disjoint Sets}]
The sum rule only holds when the sets $A_1, \ldots, A_k$ are pairwise disjoint. If a string (or
object) can belong to more than one category, simply adding the sizes leads to
\emph{over-counting}. Use the Inclusion-Exclusion Principle instead (Section~\ref{sec:ie}).
\end{warning}

% =============================================================================
\section{The Inclusion-Exclusion Principle}
\label{sec:ie}
% =============================================================================

\subsection{Two Sets}

When two sets $A_1$ and $A_2$ overlap, the sum $|A_1| + |A_2|$ counts every element in
$A_1 \cap A_2$ twice. We correct for this by subtracting the intersection once:

\begin{keyresult}[title={Key Result: Inclusion-Exclusion (Two Sets)}]
For any two sets $A_1$ and $A_2$:
\[
  |A_1 \cup A_2| = |A_1| + |A_2| - |A_1 \cap A_2|.
\]
When $A_1$ and $A_2$ are disjoint, $|A_1 \cap A_2| = 0$ and this reduces to the Sum Rule.
\end{keyresult}

\begin{example}[Bit strings (revisited)]
\label{ex:bitstringIE}
How many 8-bit binary strings either start with a ``1'' bit or end with the two bits ``00''?

\textit{Solution.} Let $A_1$ be the set of strings starting with ``1'' and $A_2$ the set
ending with ``00''. These sets are \emph{not} disjoint (a string can do both).

\begin{align*}
  |A_1| &= 2^7 = 128 \quad \text{(first bit fixed, 7 free)}, \\
  |A_2| &= 2^6 = 64 \quad \text{(last two bits fixed, 6 free)}, \\
  |A_1 \cap A_2| &= 2^5 = 32 \quad \text{(first bit and last two bits fixed, 5 free)}.
\end{align*}

By inclusion-exclusion:
\[
  |A_1 \cup A_2| = 128 + 64 - 32 = 160.
\]
\end{example}

\begin{example}[Students taking courses]
There are 1807 students in SPMS. Of these, 453 take Mathematics, 567 take Physics, and 299
take both. How many students take neither?

\textit{Solution.} Let $M$ and $P$ be the sets of students taking Maths and Physics
respectively. By inclusion-exclusion:
\[
  |M \cup P| = 453 + 567 - 299 = 721.
\]
Hence $1807 - 721 = 1086$ students take neither course.
\end{example}

\subsection{Three Sets}

\begin{theorem}[Inclusion-Exclusion for Three Sets]
\label{thm:ie3}
\[
  |A \cup B \cup C| = |A| + |B| + |C| - |A \cap B| - |A \cap C| - |B \cap C| + |A \cap B \cap C|.
\]
\end{theorem}

The pattern: add individual sets, subtract pairwise intersections, add the triple intersection.
This compensates for the layered over-counting introduced by overlapping regions in the Venn
diagram.

\begin{example}[Language courses]
1232 students take Spanish, 879 take French, 114 take Russian; 103 take both Spanish and
French, 23 take both Spanish and Russian, 14 take both French and Russian. If 2092 students
take at least one course, how many take all three?

\textit{Solution.} Let $S$, $F$, $R$ denote the three sets. By Theorem~\ref{thm:ie3}:
\[
  2092 = 1232 + 879 + 114 - 103 - 23 - 14 + |S \cap F \cap R|.
\]
Solving, $|S \cap F \cap R| = 2092 - 2085 = 7$. So 7 students take all three courses.
\end{example}

\subsection{General Inclusion-Exclusion}

\begin{theorem}[Inclusion-Exclusion Principle (General)]
\label{thm:ieg}
For sets $A_1, A_2, \ldots, A_n$:
\begin{align*}
  \left|\bigcup_{i=1}^{n} A_i\right| =
    &\sum_{i=1}^{n} |A_i|
    - \sum_{1 \le i < j \le n} |A_i \cap A_j|
    + \sum_{1 \le i < j < k \le n} |A_i \cap A_j \cap A_k|
    - \cdots \\
    &+ (-1)^{n-1} |A_1 \cap A_2 \cap \cdots \cap A_n|.
\end{align*}
Equivalently: add individual sizes, subtract all two-way intersections, add all three-way
intersections, subtract all four-way intersections, and so on, alternating signs.
\end{theorem}

\begin{proof}[Sketch of Proof]
By induction on $n$. The base case $n=1$ is trivial. For $n=2$, note that $A = (A \setminus B)
\cup (A \cap B)$ is a disjoint union, so $|A| = |A \setminus B| + |A \cap B|$. Since
$A \cup B = B \cup (A \setminus B)$ is also a disjoint union:
\[
  |A \cup B| = |B| + |A \setminus B| = |B| + |A| - |A \cap B|.
\]
For the inductive step, given the formula holds for $n$ sets, write
$\bigcup_{i=1}^{n+1} A_i = \left(\bigcup_{i=1}^{n} A_i\right) \cup A_{n+1}$
and apply the two-set formula. Using the identity
$\left(\bigcup_{i=1}^{n} A_i\right) \cap A_{n+1} = \bigcup_{i=1}^{n} (A_i \cap A_{n+1})$,
apply the inductive hypothesis to both $\bigcup_{i=1}^n A_i$ and
$\bigcup_{i=1}^n (A_i \cap A_{n+1})$ and simplify to obtain the formula for $n+1$ sets.
\end{proof}

% =============================================================================
\section*{Reading Assignment}
% =============================================================================

\begin{itemize}
  \item Rosen (7th edition), Chapter 6.1: The Basics of Counting.
  \item Rosen (7th edition), Chapter 8.5: Inclusion-Exclusion.
\end{itemize}

% =============================================================================
\chapter{Logic and Proof Techniques}
\label{chap:logic}

\section{Conditional Statements}
\label{sec:conditional}
% =============================================================================

This section covers conditional (if-then) statements, their logical properties, and the key
proof techniques that rely on them: proof by contradiction and proof by contraposition.

\subsection{Definition and Truth Table}

\begin{definition}[Conditional Statement]
\label{def:conditional}
Let $p$ and $q$ be statement variables. The \textbf{conditional} of $q$ by $p$, denoted
$p \to q$, is the statement ``If $p$ then $q$'', or equivalently ``$p$ implies $q$''. We
call $p$ the \textbf{hypothesis} (or \textbf{antecedent}) and $q$ the \textbf{conclusion}
(or \textbf{consequent}).
\end{definition}

The truth value of $p \to q$ is determined by the following truth table:

\[
\begin{array}{cc|c}
  p & q & p \to q \\ \hline
  1 & 1 & 1 \\
  1 & 0 & 0 \\
  0 & 1 & 1 \\
  0 & 0 & 1 \\
\end{array}
\]

The conditional is \emph{false only} when $p$ is true and $q$ is false. When $p$ is false,
$p \to q$ is always true, regardless of $q$. In this case we say the statement is
\textbf{vacuously true} (or \textbf{true by default}).

\begin{example}
``If I am Superman, then you are Spiderman.'' Since the hypothesis is false (no one is
Superman), the statement is vacuously true.
\end{example}

\begin{keyresult}[title={Key Result: When is $p \to q$ false?}]
$p \to q$ is \textbf{false if and only if} $p$ is true and $q$ is false. In all other
cases it is true.
\end{keyresult}

\subsection{Logical Equivalence: \texorpdfstring{$p \to q \equiv \lnot p \lor q$}{p implies q}}
\label{subsec:ptoq-equiv}

A useful rewriting of the conditional is:
\[
  p \to q \;\equiv\; \lnot p \lor q.
\]
This can be verified directly via a truth table (Tutorial exercise).

\textbf{Application.} Any disjunction of the form $\lnot p \lor q$ can be rewritten in
if-then form. For instance, ``Either you get to work on time or you are fired'' can be
expressed as ``If you do not get to work on time, then you are fired.'' (Equivalently, by
the contrapositive discussed below, this is the same as ``If you are not fired, then you
got to work on time.'')

\subsection{Negation of a Conditional}

\begin{proposition}
\label{prop:neg-conditional}
$\lnot(p \to q) \;\equiv\; p \land \lnot q$.
\end{proposition}

\begin{proof}
Using the equivalence from Section~\ref{subsec:ptoq-equiv} and De Morgan's law:
\[
  \lnot(p \to q)
  \;\equiv\; \lnot(\lnot p \lor q)
  \;\equiv\; \lnot(\lnot p) \land \lnot q
  \;\equiv\; p \land \lnot q. \qedhere
\]
\end{proof}

\begin{example}
\begin{itemize}[nosep]
  \item ``If you attend all lectures, then you will get A's for all of your courses.''\\
        \textit{Negation:} You attended all lectures and you did not get A's for all of your
        courses.
  \item ``If you are caught littering, then you will be fined.''\\
        \textit{Negation:} You are caught littering but you are not fined.
\end{itemize}
\end{example}

\begin{warning}[title={Warning: Negating an If-Then Statement}]
The negation of an if-then statement does \textbf{not} start with the word ``if''. It takes
the form $p \land \lnot q$, which is a conjunction.
\end{warning}

\subsection{Contrapositive}

\begin{definition}[Contrapositive]
\label{def:contrapositive}
The \textbf{contrapositive} of $p \to q$ is $\lnot q \to \lnot p$.
\end{definition}

A conditional and its contrapositive are logically equivalent:
\[
  p \to q \;\equiv\; \lnot q \to \lnot p.
\]

% =============================================================================
\section{Proof Techniques}
\label{sec:proofs}
% =============================================================================

\subsection{Getting Proofs Started}

Every proof of a conditional statement has a \emph{starting point} (the hypothesis) and a
\emph{conclusion to be shown}. It is good practice to write these out explicitly at the
beginning and end of the proof respectively. For example, to prove that a group of order 36
is not simple, one begins: ``Let $G$ be a group of order 36.'' and concludes:
``$\therefore G$ is not simple.''

\subsection{Proof by Contradiction}

\begin{keyresult}[title={Method: Proof by Contradiction}]
\begin{enumerate}[nosep]
  \item \textbf{Suppose the statement is false}, i.e.\ assume the negation is true.
  \item \textbf{Derive a contradiction} from this assumption using logical deduction.
  \item \textbf{Conclude} that the original statement must be true.
\end{enumerate}
Be careful when writing the negation in Step~1.
\end{keyresult}

\begin{theorem}[Theorem 4.6.1]
\label{thm:no-greatest-int}
There is no greatest integer.
\end{theorem}

\begin{proof}
Suppose not. Then there exists a greatest integer $N \in \mathbb{Z}$, so $m \le N$ for all
$m \in \mathbb{Z}$. Take $M := N + 1$. Then $M \in \mathbb{Z}$ and $M > N$, contradicting
$N$ being the greatest integer.
\end{proof}

\begin{theorem}[Theorem 4.6.2]
\label{thm:not-even-odd}
There is no integer that is both even and odd.
\end{theorem}

\begin{proof}
Suppose not. Then there exists $N \in \mathbb{Z}$ that is both even and odd. Since $N$ is
even, $N = 2k$ for some $k \in \mathbb{Z}$. Since $N$ is odd, $N = 2l + 1$ for some
$l \in \mathbb{Z}$. Hence:
\[
  2k = 2l + 1 \implies 2(k - l) = 1 \implies k - l = \tfrac{1}{2} \notin \mathbb{Z}.
\]
But $k - l \in \mathbb{Z}$ (difference of two integers), a contradiction.
\end{proof}

\begin{theorem}[Theorem 4.6.3]
\label{thm:rat-irrat-irrat}
The sum of any rational number and any irrational number is irrational.
\end{theorem}

\begin{proof}
Suppose not. Then there exist $r, s \in \mathbb{R}$ with $r \in \mathbb{Q}$,
$s \notin \mathbb{Q}$, and $r + s \in \mathbb{Q}$. Write $r = \tfrac{a}{b}$ and
$r + s = \tfrac{c}{d}$ for integers $a, b, c, d$ with $b, d \ne 0$. Then:
\[
  s = (r + s) - r = \frac{c}{d} - \frac{a}{b} = \frac{bc - ad}{bd}.
\]
Since $bc - ad \in \mathbb{Z}$ and $bd \ne 0$ (zero product rule), we get $s \in \mathbb{Q}$,
contradicting $s \notin \mathbb{Q}$.
\end{proof}

\subsection{Proof by Contraposition}

\begin{keyresult}[title={Method: Proof by Contraposition}]
To prove $\forall x \in D,\ P(x) \Rightarrow Q(x)$:
\begin{enumerate}[nosep]
  \item Rewrite as the contrapositive: $\forall x \in D,\ \lnot Q(x) \Rightarrow \lnot P(x)$.
  \item Suppose $x \in D$ with $Q(x)$ false (i.e.\ $\lnot Q(x)$ holds).
  \item Show directly that $P(x)$ is also false (i.e.\ $\lnot P(x)$ holds).
\end{enumerate}
\end{keyresult}

\begin{proposition}[Proposition 4.6.4]
\label{prop:n2even}
For all integers $n$, if $n^2$ is even then $n$ is even.
\end{proposition}

\begin{proof}[Proof (by contraposition)]
We prove the contrapositive: if $n$ is odd then $n^2$ is odd. Suppose $n$ is odd, so
$n = 2k + 1$ for some $k \in \mathbb{Z}$. Then:
\[
  n^2 = (2k+1)^2 = 4k^2 + 4k + 1 = 2(2k^2 + 2k) + 1,
\]
which is odd.
\end{proof}

\subsection{Relationship Between Contradiction and Contraposition}

The same chain of deductions used in a proof by contraposition can be recast as a proof by
contradiction. To prove $\forall x \in D,\ P(x) \Rightarrow Q(x)$ by contradiction:

\begin{enumerate}
  \item Assume $P(x)$ and $\lnot Q(x)$ simultaneously.
  \item Follow the contrapositive argument to derive $\lnot P(x)$.
  \item This gives $P(x) \land \lnot P(x)$, a contradiction ($\bot$).
\end{enumerate}

\begin{example}[Proposition 4.6.4 reproved by contradiction]
Suppose $n^2$ is even and $n$ is odd. Since $n$ is odd, write $n = 2k+1$ and compute
$n^2 = 2(2k^2 + 2k) + 1$, so $n^2$ is odd. Thus $n^2$ is both even and odd, contradicting
Theorem~\ref{thm:not-even-odd}.
\end{example}

\begin{examtip}[title={Exam Tip: Choosing Between Contradiction and Contraposition}]
Both methods are valid for conditional statements. Proof by \emph{contraposition} is often
cleaner when the negation of the conclusion has a concrete algebraic form (e.g.\ ``$n$ is
odd''). Proof by \emph{contradiction} is preferred when the statement is existential (e.g.\
``there is no greatest integer'') or when no obvious contrapositive structure exists.
\end{examtip}

%% ============================================================
%%  HANDOUT 2 --- Pigeonhole Principle, Permutations, Combinations
%% ============================================================

\newpage
\section{The Pigeonhole Principle}
\label{sec:pigeonhole}

\subsection{Motivating Example}

\begin{example}
Suppose a flock of 20 pigeons flies into 19 pigeonholes to roost. Prove that at least one
pigeonhole must contain at least two pigeons.
\end{example}

\begin{proof}
Suppose for contradiction that every pigeonhole has at most one pigeon. Then across all 19
pigeonholes there are at most 19 pigeons, which contradicts the fact that 20 pigeons were placed.
\end{proof}

\subsection{The Pigeonhole Principle}

\begin{theorem}[Pigeonhole Principle (Section 6.2)]
\label{thm:pigeonhole}
Let $k$ be a positive integer. If $k+1$ or more objects are placed into $k$ boxes, then
there is at least one box containing two or more objects.
\end{theorem}

\begin{proof}
We prove the contrapositive: if no box contains more than one object, then fewer than $k+1$
objects are placed into the $k$ boxes. Indeed, if every box holds at most one object, the total
number of objects is at most $k < k+1$.
\end{proof}

\begin{keyresult}[title={Key Result: Pigeonhole Principle}]
If $n$ objects are placed into $k$ boxes with $n > k$, then some box contains $\geq 2$ objects.
The \emph{objects} are the ``pigeons''; the \emph{boxes} (categories) are the ``pigeonholes''.
Identifying the right assignment is the key step in applying this principle.
\end{keyresult}

\subsection{Application: Multiples with Only 0s and 1s}

\begin{example}[Example 4, Section 6.2]
Show that for every positive integer $n$, there is a multiple of $n$ whose decimal expansion
consists entirely of $0$s and $1$s.
\end{example}

\begin{proof}
Consider the set of $n+1$ integers
\[
  X = \{1,\; 11,\; 111,\; \ldots,\; \underbrace{11\cdots 1}_{n+1}\}.
\]
There are $n$ possible remainders when an integer is divided by $n$ (namely $0, 1, \ldots, n-1$).
Since $|X| = n+1 > n$, by the Pigeonhole Principle two distinct elements $x > x'$ in $X$ share the
same remainder when divided by $n$. Therefore $x - x'$ is a multiple of $n$, and its decimal
expansion consists entirely of $0$s and $1$s (e.g.\ $111\,100\cdots 0 = 1111\cdots 1 - 11\cdots 1$).
\end{proof}

\begin{examtip}[title={Exam Tip: Identifying Pigeons and Pigeonholes}]
Always state explicitly:
\begin{itemize}
  \item What are the \emph{pigeons} (objects)?
  \item What are the \emph{pigeonholes} (boxes / categories)?
  \item Why does $|\text{pigeons}| > |\text{pigeonholes}|$?
\end{itemize}
The conclusion then follows automatically from the theorem.
\end{examtip}

\subsection{Recall: Floor and Ceiling Functions}

\begin{definition}[Floor and Ceiling]
\label{def:floor-ceiling}
For any real number $x$:
\begin{itemize}
  \item The \textbf{floor} of $x$, written $\lfloor x \rfloor$, is the unique integer $n$ such
        that $n \leq x < n+1$.
  \item The \textbf{ceiling} of $x$, written $\lceil x \rceil$, is the unique integer $n$ such
        that $n-1 < x \leq n$.
\end{itemize}
Symbolically, $\lfloor x \rfloor = n \iff n \leq x < n+1$ and
$\lceil x \rceil = n \iff n-1 < x \leq n$.
\end{definition}

Examples: $\lfloor 4.5 \rfloor = 4$, $\lfloor -1.73 \rfloor = -2$, $\lfloor -5 \rfloor = -5$;
$\lceil 4.5 \rceil = 5$, $\lceil -4.5 \rceil = -4$, $\lceil -4 \rceil = -4$.

\subsection{Generalised Pigeonhole Principle}

\begin{theorem}[Generalised Pigeonhole Principle (Section 6.2)]
\label{thm:gen-pigeonhole}
If $N$ objects are placed into $k$ boxes, then at least one box contains at least
$\lceil N/k \rceil$ objects.
\end{theorem}

\begin{proof}
Suppose for contradiction that every box contains fewer than $\lceil N/k \rceil$ objects, i.e.\
each box holds at most $\lceil N/k \rceil - 1$ objects. Then the total number of objects is at most
\[
  k\!\left(\lceil N/k \rceil - 1\right)
  < k\!\left(\frac{N}{k} + 1 - 1\right)
  = N,
\]
where we used $\lceil N/k \rceil < N/k + 1$. This contradicts the existence of $N$ objects.
\end{proof}

\begin{example}[Hair Coincidence]
Singapore has more than $5\,000\,000$ non-bald people. Assume every person has fewer than
$200\,000$ hairs. Show that at least 26 people have the same number of hairs.
\end{example}

\begin{proof}
Let $N \geq 5\,000\,001$ (people) and $k \leq 199\,999$ (possible hair counts $0,1,\ldots,199\,999$).
By the Generalised Pigeonhole Principle some hair-count category contains at least
\[
  \left\lceil \frac{5\,000\,001}{199\,999} \right\rceil = 26
\]
people.
\end{proof}

\begin{examtip}[title={Exam Tip: Generalised Pigeonhole Ceiling Computation}]
For 20 objects distributed among 3 boxes (e.g.\ 20 students choosing from 3 burger types):
$\lceil 20/3 \rceil = 7$. So at least 7 students chose the same type. This is \emph{always}
true — it is not possible to have every type chosen by at most 6 (that would give at most
$6 \times 3 = 18 < 20$).
\end{examtip}

%% ──────────────────────────────────────────────────────────────
\chapter{Counting Methods}
\label{chap:counting-methods}

\section{Permutations}
\label{sec:permutations}

\subsection{Definitions}

\begin{definition}[Permutation and $r$-Permutation]
\label{def:r-perm}
A \textbf{permutation} of a set $S$ of distinct elements is an \emph{ordered} arrangement of
all elements of $S$.

An \textbf{$r$-permutation} of $S$ (with $|S| = n$, $r \leq n$) is an \emph{ordered} selection
of $r$ distinct elements from $S$.

The number of $r$-permutations of an $n$-element set is denoted $P(n,r)$, also written
${}^n P_r$ or $nP_r$.
\end{definition}

\begin{example}
Let $S = \{a, b, c\}$. Then $(b, c, a)$ is a permutation of $S$. The $2$-permutations of $S$
are $(a,b),(b,a),(a,c),(c,a),(b,c),(c,b)$, so $P(3,2) = 6$.
\end{example}

\subsection{Formula for $P(n,r)$}

\begin{theorem}[Section 6.3]
\label{thm:pnr}
For positive integers $n$ and $r$ with $1 \leq r \leq n$,
\[
  P(n,r) = n(n-1)(n-2)\cdots(n-r+1).
\]
\end{theorem}

\begin{proof}
To build an $r$-permutation, choose elements one position at a time. The first position has $n$
choices, the second $n-1$, and so on; the $r$-th position has $n-(r-1) = n-r+1$ choices.
By the Product Rule (see Section~\ref{sec:conditional}),
$P(n,r) = n(n-1)\cdots(n-r+1)$.
\end{proof}

\begin{corollary}[Section 6.3]
\label{cor:pnr-factorial}
For integers $n, r$ with $0 \leq r \leq n$,
\[
  P(n,r) = \frac{n!}{(n-r)!}.
\]
\end{corollary}

\textbf{Special cases:}
\begin{itemize}
  \item $P(n,n) = n!$ \quad (all permutations of $n$ elements; recall $0! = 1$).
  \item $P(n,1) = n$.
  \item $P(n,0) = 1$ (exactly one way to arrange zero elements).
\end{itemize}

Useful factorials to memorise: $1!=1,\; 2!=2,\; 3!=6,\; 4!=24,\; 5!=120,\; 6!=720$.

\begin{keyresult}[title={Key Result: $P(n,r) = n!/(n-r)!$}]
Order matters in a permutation. Use $P(n,r)$ whenever you are selecting $r$ items from $n$
\emph{and the order of selection is significant}.
\end{keyresult}

\subsection{Examples}

\begin{example}[COMPUTER]
\begin{itemize}
  \item[(i)] The letters of \textsc{COMPUTER} can be arranged in $P(8,8) = 8! = 40\,320$ ways.
  \item[(ii)] If \texttt{CO} must appear as a fixed unit, treat it as a single character. There
              are then 7 distinct characters, giving $P(7,7) = 7! = 5\,040$ arrangements.
\end{itemize}
\end{example}

\begin{example}[Prize Winners]
The number of ways to award first, second, and third prizes from 200 students is
\[
  P(200, 3) = 200 \times 199 \times 198 = 7\,880\,400.
\]
\end{example}

\begin{example}[EQUATION -- $r$-permutations with a constraint]
\begin{itemize}
  \item[(i)] The number of 3-letter arrangements from the 8 distinct letters of \textsc{EQUATION}
             is $P(8,3) = 8 \times 7 \times 6 = 336$.
  \item[(ii)] The number of 4-letter arrangements containing \texttt{QU} together (in that order):
              treat \texttt{QU} as one unit, leaving 2 more letters chosen from the remaining 6.
              Place the \texttt{QU}-block in any of the 3 slots, and fill the other 2 from 6
              remaining letters. Two sub-cases --- \texttt{QU} first or last --- each give $6 \times 5 = 30$,
              for a total of $30 + 30 = \mathbf{90}$.
\end{itemize}
\end{example}

%% ──────────────────────────────────────────────────────────────
\section{Combinations}
\label{sec:combinations}

\subsection{Definition}

\begin{definition}[$r$-Combination]
\label{def:r-comb}
An \textbf{$r$-combination} of a set $S$ with $n$ elements is an \emph{unordered} selection of
$r$ elements from $S$ (i.e.\ a subset of $S$ of size $r$), where $r \leq n$.

The number of $r$-combinations of an $n$-element set is denoted $C(n,r)$, $\binom{n}{r}$,
${}^n C_r$, or $nC_r$, and is called a \textbf{binomial coefficient}, read ``$n$ choose $r$''.
\end{definition}

\begin{warning}[title={Warning: Permutation vs.\ Combination}]
\begin{itemize}
  \item \textbf{Permutation (ordered):} $(a,b)$ and $(b,a)$ are \emph{different}.
  \item \textbf{Combination (unordered):} $\{a,b\}$ and $\{b,a\}$ are the \emph{same}.
\end{itemize}
Ask: ``Does the \emph{order} of selection matter?'' If yes, use $P(n,r)$. If no, use $C(n,r)$.
\end{warning}

\subsection{Relation Between $P(n,r)$ and $C(n,r)$}

Building an $r$-permutation is a two-step process:
\begin{enumerate}
  \item Choose an $r$-element subset (combination): $C(n,r)$ ways.
  \item Order the chosen subset: $P(r,r) = r!$ ways.
\end{enumerate}
By the Product Rule, $P(n,r) = C(n,r) \cdot r!$, so
\[
  C(n,r) = \frac{P(n,r)}{r!} = \frac{n!}{r!\,(n-r)!}.
\]

\subsection{Formula for $C(n,r)$}

\begin{theorem}[Section 6.3]
\label{thm:cnr}
For integers $r, n$ with $0 \leq r \leq n$,
\[
  C(n,r) = \binom{n}{r} = \frac{n!}{r!\,(n-r)!}.
\]
\end{theorem}

\textbf{Special cases:}
\begin{itemize}
  \item $C(n,n) = 1$ (only one way to choose all $n$ elements; uses $0! = 1$).
  \item $C(n,1) = n$.
  \item $C(n, n-1) = n$.
  \item $C(n,0) = 1$.
\end{itemize}

\begin{keyresult}[title={Key Result: $C(n,r) = n!/[r!(n-r)!]$}]
Use $C(n,r)$ when selecting $r$ items from $n$ with \emph{no regard to order}.
The formula satisfies $C(n,r) = C(n, n-r)$ — choosing $r$ items to \emph{include} is equivalent
to choosing $n-r$ items to \emph{exclude}.
\end{keyresult}

\subsection{Symmetry of Binomial Coefficients}

\begin{corollary}[Section 6.3]
\label{cor:symmetry}
For nonnegative integers $r \leq n$,
\[
  C(n, r) = C(n, n-r).
\]
\end{corollary}

\begin{proof}[Proof 1 (Algebraic)]
$C(n,r) = \dfrac{n!}{r!(n-r)!} = \dfrac{n!}{(n-r)!\,[n-(n-r)]!} = C(n, n-r)$.
\end{proof}

\begin{proof}[Proof 2 (Combinatorial)]
Every $r$-element subset $A$ of an $n$-element set $S$ corresponds bijectively to its complement
$\overline{A}$, which has $n-r$ elements. Thus there are equally many $r$-subsets and
$(n-r)$-subsets.
\end{proof}

\subsection{Examples}

\begin{example}[0-1 Bit Strings, Example 14 Section 6.3]
The number of 0-1 bit strings of length $n$ containing exactly $r$ ones equals $C(n,r)$, since
we are choosing which $r$ of the $n$ positions to fill with $1$.
\end{example}

\begin{example}[Forming a Team, Example 15 Section 6.3]
There are 9 boys and 11 girls. The number of ways to choose a team of 3 boys and 4 girls is
\[
  C(9,3) \cdot C(11,4) = 84 \times 330 = 27\,720.
\]
\end{example}

\begin{example}[A and B Must Work Together]
From 12 people (including A and B who must both be on the team or neither), the number of
5-person teams is:
\begin{itemize}
  \item Both A and B included: choose 3 more from remaining 10: $C(10,3) = 120$.
  \item Neither A nor B included: choose 5 from remaining 10: $C(10,5) = 252$.
\end{itemize}
By the Sum Rule: $120 + 252 = \mathbf{372}$.
\end{example}

\begin{example}[A and B Cannot Work Together]
If A and B cannot both be on the team, the three cases are:
\begin{itemize}
  \item A only (B excluded): choose 4 more from remaining 10: $C(10,4) = 210$.
  \item B only (A excluded): $C(10,4) = 210$.
  \item Neither A nor B: $C(10,5) = 252$.
\end{itemize}
Total: $210 + 210 + 252 = \mathbf{672}$.
\end{example}

\begin{example}[Powerball Lottery]
In Powerball, 5 white balls are chosen from 69 (distinct, unordered) and 1 red ``Powerball''
from 26. The total number of combinations is
\[
  C(69, 5) \times 26 = 11\,238\,513 \times 26 = 292\,201\,338.
\]
\end{example}

\begin{example}[Poker Hands]
\begin{itemize}
  \item[(i)] Five-card hands from 52 cards: $C(52,5) = \dfrac{52!}{5!\,47!} = 2\,598\,960$.
  \item[(ii)] The number of ways to select 47 cards is $C(52,47) = C(52,5) = 2\,598\,960$,
              by the symmetry $C(n,r) = C(n,n-r)$.
\end{itemize}
\end{example}

%% ──────────────────────────────────────────────────────────────
\section{Combinatorial Proofs}
\label{sec:combproof}

\begin{definition}[Combinatorial Proof]
\label{def:combproof}
A \textbf{combinatorial proof} of an identity is a proof showing that both sides count the same
collection of objects, but in different ways (also called a \emph{double-counting argument}).
\end{definition}

\subsection{Pascal's Identity}

\begin{theorem}[Pascal's Identity]
\label{thm:pascal}
For all integers $n \geq k \geq 1$,
\[
  \binom{n}{k} = \binom{n-1}{k} + \binom{n-1}{k-1}.
\]
\end{theorem}

\begin{proof}[Proof 1 (Algebraic)]
\begin{align*}
\binom{n-1}{k-1} + \binom{n-1}{k}
  &= \frac{(n-1)!}{(k-1)!\,(n-k)!} + \frac{(n-1)!}{k!\,(n-1-k)!} \\
  &= \frac{(n-1)!\,k}{k!\,(n-k)!} + \frac{(n-1)!\,(n-k)}{k!\,(n-k)!} \\
  &= \frac{(n-1)!\,(k + n - k)}{k!\,(n-k)!}
   = \frac{n!}{k!\,(n-k)!} = \binom{n}{k}.
\end{align*}
\end{proof}

\begin{proof}[Proof 2 (Combinatorial)]
We count $k$-subsets of $S = \{1, 2, \ldots, n\}$ in two ways. Fix element $1$.
\begin{itemize}
  \item \textbf{Case 1 — element $1$ is in the subset:} choose $k-1$ more from the remaining
        $n-1$ elements: $\binom{n-1}{k-1}$ ways.
  \item \textbf{Case 2 — element $1$ is not in the subset:} choose all $k$ from the remaining
        $n-1$ elements: $\binom{n-1}{k}$ ways.
\end{itemize}
Since these cases partition all $k$-subsets: $\binom{n}{k} = \binom{n-1}{k-1} + \binom{n-1}{k}$.
\end{proof}

\subsection{The 6-Person Party Problem}

\begin{example}
Among any party of 6 people, there exist 3 people who all know one another or 3 who are all
strangers to one another.
\end{example}

\begin{proof}[Sketch]
Fix one person $A$. There are 5 edges from $A$ to others, each coloured ``know'' or ``stranger''.
By the Pigeonhole Principle (Theorem~\ref{thm:pigeonhole}), at least $\lceil 5/2 \rceil = 3$ edges
share the same colour — say $A$ knows $B, C, D$.
\begin{itemize}
  \item If any pair among $\{B, C, D\}$ know each other (say $B$ knows $C$), then $A, B, C$
        all know each other.
  \item If no pair among $\{B, C, D\}$ know each other, then $B, C, D$ are all mutual strangers.
\end{itemize}
In either case, a monochromatic triple exists.
\end{proof}

\subsection{Social Connections: Two People with Equal Degree}

\begin{example}
Among $n \geq 2$ people, prove there always exist two people who know the same number of others.
\end{example}

\begin{proof}
Each person's acquaintance count lies in $\{0, 1, \ldots, n-1\}$, giving $n$ possible values.
However, $0$ and $n-1$ cannot both occur simultaneously (if someone knows everyone, no one is a
stranger to all). Hence the $n$ counts in fact lie in at most $n-1$ distinct values. By the
Pigeonhole Principle (Theorem~\ref{thm:pigeonhole}), two people share the same count.
\end{proof}

\subsection{Bijective Combinatorial Proof}

\begin{example}[Equal Partition Counts]
Let $n$ be a positive integer. Prove that the number of ways to write $n = a + b$ with integers
$a \geq b \geq 0$ equals the number of ways to write $n = c + 2d$ with nonnegative integers
$c, d \geq 0$.
\end{example}

\begin{proof}
Let $A = \{(a,b) \in \mathbb{Z}^2 : n = a+b,\; a \geq b \geq 0\}$ and
$B = \{(c,d) \in \mathbb{Z}^2 : n = c+2d,\; c, d \geq 0\}$.
Define $\varphi : A \to B$ by $\varphi(a, b) = (a-b, b)$.

\textit{Well-defined:} $a - b \geq 0$ since $a \geq b$, and $(a-b) + 2b = a+b = n$. \checkmark

\textit{Injective:} If $\varphi(a,b) = \varphi(a', b')$ then $(a-b, b) = (a'-b', b')$, so $b = b'$
and hence $a = a'$. \checkmark

\textit{Surjective:} Given $(c, d) \in B$, note $\varphi(c+d, d) = (c, d)$ and
$(c+d) + d = c + 2d = n$, $c+d \geq d \geq 0$. \checkmark

Thus $\varphi$ is a bijection, so $|A| = |B|$.
\end{proof}

%% ============================================================
%%  Handout 3: Binomial Theorem, Generalised Permutations and Combinations
%% ============================================================

\section{The Binomial Theorem}
\label{sec:binomial}

\subsection{Binomial Coefficients}

Recall that $C(n,k)$ (also written $\binom{n}{k}$) counts the number of $k$-element subsets of an
$n$-element set.  These are called \emph{binomial coefficients} because they appear as the
coefficients in the expansion of $(x+y)^n$.

To see why, consider expanding $(x+y)^n = (x+y)(x+y)\cdots(x+y)$ ($n$ factors).  Each term of the
expansion is formed by choosing either $x$ or $y$ from every factor.  A term $x^{n-i}y^{i}$ arises
exactly when we choose $y$ from $i$ of the $n$ factors (and $x$ from the remaining $n-i$); there are
$C(n,i)$ ways to make that choice.

\begin{theorem}[The Binomial Theorem]
\label{thm:binomial}
Let $x$ and $y$ be variables and let $n$ be a nonnegative integer.  Then
\[
  (x+y)^n = \sum_{i=0}^{n} C(n,i)\,x^{n-i}y^{i}.
\]
\end{theorem}

\begin{proof}[Combinatorial proof]
Expanding $(x+y)(x+y)\cdots(x+y)$, a term of the form $x^{n-i}y^{i}$ is obtained by choosing
$x$ from exactly $n-i$ of the $n$ factors (equivalently, choosing $y$ from the remaining $i$
factors).  The number of such choices is $C(n,n-i) = C(n,i)$.  Summing over $i$ gives the result.
\end{proof}

\begin{keyresult}[title={Key Result: Binomial Theorem}]
\[
  (x+y)^n = \sum_{i=0}^{n} \binom{n}{i} x^{n-i} y^{i}.
\]
The coefficient of $x^{n-i}y^{i}$ is $C(n,i)$, the number of ways to choose $i$ factors
from which to take $y$.
\end{keyresult}

\begin{example}
Expand $(x+y)^4$.

\[
  (x+y)^4 = C(4,0)x^4 + C(4,1)x^3y + C(4,2)x^2y^2 + C(4,3)xy^3 + C(4,4)y^4
           = x^4 + 4x^3y + 6x^2y^2 + 4xy^3 + y^4.
\]
\end{example}

\begin{example}
What is the coefficient of $x^{12}y^{13}$ in the expansion of $(x+y)^{25}$?

By the Binomial Theorem (Theorem~\ref{thm:binomial}), the coefficient is
\[
  C(25,13) = \frac{25!}{13!\,12!} = 5{,}200{,}300.
\]
\end{example}

\begin{example}
\label{ex:binomial-scaled}
What is the coefficient of $x^{12}y^{13}$ in the expansion of $(2x - 3y)^{25}$?

Write $(2x-3y)^{25} = [(2x)+(-3y)]^{25} = \sum_{i=0}^{25} C(25,i)(2x)^{25-i}(-3y)^{i}$.
The term with $i=13$ gives $x^{12}y^{13}$, so the coefficient is
\[
  C(25,13)\,2^{12}\,(-3)^{13} = -\frac{25!}{13!\,12!}\,2^{12}\,3^{13}.
\]
\end{example}

\begin{examtip}[title={Exam Tip: Scaled Binomial}]
When the expression is $(ax + by)^n$, substitute $X = ax$ and $Y = by$ to apply the
standard Binomial Theorem, then match the desired monomial.  Remember to include the
scaling factors $a^{n-i}$ and $b^i$ in the coefficient.
\end{examtip}

\subsection{Pascal's Triangle}

The binomial coefficients $C(n,i)$ for successive values of $n$ form Pascal's Triangle:

\[
\begin{array}{ccccccccc}
  n=0: & & & & 1 \\
  n=1: & & & 1 & & 1 \\
  n=2: & & 1 & & 2 & & 1 \\
  n=3: & 1 & & 3 & & 3 & & 1 \\
  n=4: & 1 & & 4 & & 6 & & 4 & & 1
\end{array}
\]

Each interior entry is the sum of the two entries directly above it, which is precisely
Pascal's Identity (Theorem~\ref{thm:pascal}).

\subsection{Corollaries of the Binomial Theorem}

\begin{corollary}
\label{cor:sum2n}
For every nonnegative integer $n$,
\[
  \sum_{i=0}^{n} C(n,i) = 2^n.
\]
\end{corollary}

\begin{proof}
Set $x = y = 1$ in the Binomial Theorem:
\[
  2^n = (1+1)^n = \sum_{i=0}^{n} C(n,i)\,1^{n-i}\,1^{i} = \sum_{i=0}^{n} C(n,i). \qedhere
\]
\end{proof}

\begin{proof}[Combinatorial proof]
Let $S$ be a set with $n$ elements.  The right-hand side $2^n$ counts all subsets of $S$
(one per binary string of length $n$).  The left-hand side counts the same subsets grouped
by size: $C(n,0)$ subsets of size $0$, $C(n,1)$ of size $1$, \dots, $C(n,n)$ of size $n$.
Both sides count the same collection, so they are equal.
\end{proof}

\begin{keyresult}[title={Key Result: Sum of Binomial Coefficients}]
The $n$-th row of Pascal's Triangle sums to $2^n$.  Equivalently, an $n$-element set has
exactly $2^n$ subsets in total.
\end{keyresult}

%% ============================================================
\section{Generalised Permutations and Combinations}
\label{sec:genpermcomb}

Real counting problems often involve elements placed in a cycle, elements used repeatedly, or
indistinguishable elements.  The key strategies are:
\begin{itemize}
  \item View the problem from a different perspective.
  \item Apply or combine the appropriate counting rules.
  \item Use bijective functions to reduce to a known case.
  \item Never overcount or undercount.
\end{itemize}

\subsection{The Division Rule}
\label{subsec:divisionrule}

\begin{definition}[$k$-to-$1$ function]
A function $f: X \to Y$ is \emph{$k$-to-$1$} if for every $y \in Y$, the number of
elements $x \in X$ with $f(x) = y$ is exactly $k$.
\end{definition}

\begin{theorem}[Division Rule]
\label{thm:division}
If $f : X \to Y$ is $k$-to-$1$, then $|Y| = \dfrac{|X|}{k}$.
\end{theorem}

The Division Rule is used when we know $|X|$ and can construct a $k$-to-$1$ map from $X$ to
the set $Y$ we want to count.

\subsection{Permutations Around a Cycle}
\label{subsec:cycleperm}

When objects are arranged around a circle, rotations of the same arrangement are considered
identical.

\begin{example}
How many ways are there to seat $3$ people around a round table?

A linear arrangement of $3$ people gives $P(3,3) = 6$ orderings.  However, the $3$ cyclic
rotations $(A,B,C)$, $(B,C,A)$, $(C,A,B)$ all represent the same circular seating.  Dividing by
$3$ gives $6/3 = 2$ distinct circular arrangements.  (The two arrangements correspond to
``$B$ is to $A$'s left'' and ``$C$ is to $A$'s left''.)
\end{example}

\begin{theorem}[Circular $r$-Permutations]
\label{thm:cycleperm}
The number of $r$-permutations of a set with $n$ elements arranged around a cycle is
\[
  \frac{1}{r}\cdot P(n,r) = \frac{n!}{r(n-r)!}.
\]
\end{theorem}

\begin{proof}
Let $P$ be the set of all (linear) $r$-permutations of $n$ elements, so $|P| = P(n,r)$.
Let $C$ be the set of all circular $r$-permutations.  Define $f : P \to C$ by placing the
elements of each linear permutation $(p_1, p_2, \ldots, p_r)$ around a cycle in order.

The map $f$ is $r$-to-$1$ because exactly $r$ linear permutations map to the same circular
permutation: the $r$ cyclic rotations
\[
  (a_1,a_2,\ldots,a_r),\quad (a_2,a_3,\ldots,a_r,a_1),\quad \ldots,\quad (a_r,a_1,\ldots,a_{r-1})
\]
all produce the same circular arrangement.  By the Division Rule (Theorem~\ref{thm:division}),
\[
  |C| = \frac{P(n,r)}{r} = \frac{n!}{r(n-r)!}. \qedhere
\]
\end{proof}

\begin{keyresult}[title={Key Result: Circular Permutations}]
Arrange $n$ distinct objects in a circle: $(n-1)!$ ways (fix one object and permute the rest).
More generally, choose and arrange $r$ objects from $n$ in a circle: $\dfrac{P(n,r)}{r}$ ways.
\end{keyresult}

\subsection{Permutations with Repetition}

\begin{theorem}
\label{thm:permrep}
The number of $r$-permutations of a set of $n$ elements \emph{with repetition allowed} is $n^r$.
\end{theorem}

\begin{proof}
Each of the $r$ positions may be filled by any of the $n$ elements independently.  By the
Product Rule, the total is $n \cdot n \cdots n = n^r$.
\end{proof}

\begin{example}
The number of strings of length $r$ over the $26$-letter English alphabet is $26^r$.
\end{example}

\subsection{Combinations with Repetition}

\begin{theorem}
\label{thm:combrep}
The number of $r$-combinations from a set of $n$ elements \emph{with repetition allowed} is
\[
  C(n+r-1,\,r) = C(n+r-1,\,n-1).
\]
\end{theorem}

\begin{proof}
We establish a bijection between $r$-combinations with repetition and binary strings of length
$n+r-1$ containing exactly $r$ ones (``balls'') and $n-1$ zeros (``dividers'').

Represent the $n$ element types as $n$ bins separated by $n-1$ dividers.  A choice of $r$
items (with repetition) from the $n$ types is encoded by placing $r$ balls into the bins: the
number of balls in bin $i$ equals how many times type $i$ was chosen.  This bijection shows
that the count equals the number of ways to choose $r$ positions for the balls among
$n+r-1$ positions, which is $C(n+r-1,r) = C(n+r-1,n-1)$.
\end{proof}

\begin{example}[Fruit selection]
How many ways are there to select $3$ pieces of fruit from a bowl containing apples, oranges,
and pears if order does not matter and repetition is allowed?

This is a $3$-combination with repetition from $n = 3$ types:
\[
  C(3+3-1,\,3) = C(5,3) = 10.
\]
\end{example}

\begin{example}[Cookies]
A cookie shop sells $4$ kinds of cookies.  How many ways can $6$ cookies be chosen (order
does not matter, repetition allowed)?

Here $n=4$, $r=6$:
\[
  C(4+6-1,\,6) = C(9,6) = 84.
\]
\end{example}

\begin{example}[Integer solutions]
\label{ex:eqnsols}
How many solutions does $x_1 + x_2 + x_3 = 11$ have with $x_1, x_2, x_3 \geq 0$ integers?

Each solution corresponds to distributing $11$ identical balls into $3$ labelled bins, i.e.,
an $11$-combination with repetition from $3$ types:
\[
  C(3+11-1,\,11) = C(13,11) = C(13,2) = 78.
\]
\end{example}

\begin{example}[Integer solutions with lower bounds]
How many solutions does $x_1 + x_2 + x_3 = 11$ have with $x_1 \geq 1$, $x_2 \geq 2$,
$x_3 \geq 3$?

Substitute $y_i = x_i - (\text{lower bound})$: let $x_1 = 1+y_1$, $x_2 = 2+y_2$,
$x_3 = 3+y_3$.  The equation becomes $y_1 + y_2 + y_3 = 5$ with $y_i \geq 0$, giving
\[
  C(3+5-1,\,5) = C(7,5) = C(7,2) = 21.
\]
\end{example}

\begin{examtip}[title={Exam Tip: Lower-Bound Substitution}]
To handle $x_i \geq a_i > 0$, substitute $y_i = x_i - a_i$ to get $y_i \geq 0$ and reduce
the right-hand side by $\sum a_i$.  Then apply Theorem~\ref{thm:combrep} to the new equation.
\end{examtip}

\subsection{Permutations with Indistinguishable Objects}

\begin{theorem}
\label{thm:indistperm}
The number of distinct permutations of $n$ objects where there are $n_1$ indistinguishable
objects of type $1$, $n_2$ of type $2$, \ldots, $n_k$ of type $k$ (with $n_1 + \cdots + n_k = n$) is
\[
  \frac{n!}{n_1!\,n_2!\cdots n_k!}.
\]
\end{theorem}

\begin{proof}
Let $P$ be the set of all permutations treating every object as distinguishable ($|P| = n!$).
Define $f: P \to C$ (where $C$ is the set of permutations of the indistinguishable objects) by
ignoring the labels.  Two distinguishable permutations map to the same indistinguishable
permutation if and only if they differ only in the internal ordering of the objects within each
type.  The number of such orderings is $n_1!\,n_2!\cdots n_k!$, so $f$ is
$(n_1!\,n_2!\cdots n_k!)$-to-$1$.  By the Division Rule,
\[
  |C| = \frac{n!}{n_1!\,n_2!\cdots n_k!}. \qedhere
\]
\end{proof}

\begin{example}[SUCCESS]
How many strings can be made by reordering the letters of the word ``SUCCESS''?

The letters are: $3$ S's, $2$ C's, $1$ U, $1$ E (total $7$).  By Theorem~\ref{thm:indistperm},
\[
  \frac{7!}{3!\,2!\,1!\,1!} = \frac{5040}{12} = 420.
\]
\end{example}

\begin{example}[MISSISSIPPI]
How many strings can be formed from the letters of ``MISSISSIPPI''?

The letters are: $4$ I's, $4$ S's, $2$ P's, $1$ M (total $11$):
\[
  \frac{11!}{4!\,4!\,2!\,1!} = 34{,}650.
\]
\end{example}

\subsection{Distributing Objects into Boxes}

The four distribution problems form a useful taxonomy.

\begin{theorem}[Distinguishable objects into distinguishable boxes]
\label{thm:distdist}
The number of ways to distribute $n$ distinguishable objects into $k$ distinguishable boxes
so that box $i$ receives $n_i$ objects ($n_1 + \cdots + n_k = n$) is
\[
  \frac{n!}{n_1!\,n_2!\cdots n_k!}.
\]
\end{theorem}

This is the same formula as Theorem~\ref{thm:indistperm}; a bijection between the two settings
can be constructed by associating position $j$ in a string with box $j$.

\medskip\noindent
\textbf{Indistinguishable objects into distinguishable boxes.}
Placing $r$ identical balls into $n$ labelled bins is equivalent to choosing an $r$-combination
with repetition from $n$ types (Theorem~\ref{thm:combrep}), giving $C(n+r-1,r)$.

\begin{example}
The number of ways to distribute $5$ indistinguishable objects into $3$ distinguishable boxes is
$C(3+5-1,5) = C(7,2) = 21$.
\end{example}

\medskip\noindent
\textbf{Distinguishable objects into indistinguishable boxes.}
There is no simple closed form.  One must sum Stirling numbers of the second kind
$S(n,j)$ over $j \leq k$.  For small cases, break into sub-cases by how many boxes are
non-empty and divide by the number of permutations of identical boxes.

\begin{example}
The number of ways to distribute $5$ distinguishable objects into $3$ indistinguishable boxes
is $S(5,1)+S(5,2)+S(5,3) = 1 + 15 + 25 = 41$.
\end{example}

\medskip\noindent
\textbf{Indistinguishable objects into indistinguishable boxes.}
This equals the number of partitions of $n$ into at most $k$ positive parts, denoted $p_k(n)$.
There is no elementary closed form in general.

\begin{example}
$p_3(5)$: the partitions of $5$ into at most $3$ parts are
$5$; $4+1$; $3+2$; $3+1+1$; $2+2+1$.
So $p_3(5) = 5$.
\end{example}

\subsection{Summary Table}

\begin{center}
\renewcommand{\arraystretch}{1.6}
\begin{tabular}{lcc}
\hline
& \textbf{$r$-permutation} & \textbf{$r$-combination} \\
\hline
Without repetition & $P(n,r) = \dfrac{n!}{(n-r)!}$ & $C(n,r) = \dfrac{n!}{(n-r)!\,r!}$ \\[6pt]
With repetition    & $n^r$ & $C(n+r-1,r)$ \\[6pt]
Indistinguishable objects & $\dfrac{n!}{n_1!\,n_2!\cdots n_k!}$ & --- \\[6pt]
Around a cycle     & $\dfrac{P(n,r)}{r} = \dfrac{n!}{r(n-r)!}$ & --- \\
\hline
\end{tabular}
\end{center}

\begin{warning}[title={Warning: Circular vs.\ Linear}]
For $n$ objects in a \emph{line} the count is $P(n,n) = n!$; around a \emph{circle} it is
$(n-1)!$.  The division by $n$ corrects for the $n$ rotations that yield the same circular
arrangement.
\end{warning}

\section*{Reading Assignment (Handout 3)}
Rosen (7th edition): Chapter~6.4 (Binomial Coefficients and Identities);
Chapter~6.5 (Generalised Permutations and Combinations).

%% ============================================================
%%  Handout 4 – Recurrence Relations: Introduction and Applications
%% ============================================================
\chapter{Recurrence Relations: Foundations}
\label{chap:recurrence-foundations}

\section{Introduction to Recurrence Relations}
\label{sec:recurrence-intro}

A \textbf{recurrence relation} expresses the $n$-th term of a sequence in terms of one or more
earlier terms.  The core idea is:

\begin{itemize}
  \item For a problem of size $n$, relate its solution to smaller sub-problems of size
        $n-1$, $n-2$, and so on.
  \item Combine this recurrence with one or more \emph{initial conditions} to fully specify
        the sequence.
\end{itemize}

Together, an initial condition and a recurrence uniquely determine every term of the sequence.
Once a closed-form expression is conjectured, it is typically verified by induction.

\subsection{Tower of Hanoi}
\label{subsec:hanoi}

\begin{example}[Tower of Hanoi]
\label{ex:hanoi}
What is the minimum number of moves needed to transfer $n$ disks from one peg to another,
subject to the rules that (1) only one disk may be moved at a time, and (2) a larger disk
may never rest on top of a smaller one?
\end{example}

Let $H_n$ denote the minimum number of moves for $n$ disks.  Computing small cases directly
gives $H_1 = 1$, $H_2 = 3$, $H_3 = 7$.

\textbf{Recurrence.}  To move $n$ disks from peg A to peg C:
\begin{enumerate}
  \item Move the top $n-1$ disks from A to B \hfill ($H_{n-1}$ moves).
  \item Move the largest disk from A to C \hfill (1 move).
  \item Move the $n-1$ disks from B to C \hfill ($H_{n-1}$ moves).
\end{enumerate}

Hence
\[
  H_n = 2H_{n-1} + 1, \qquad H_1 = 1.
\]

\textbf{Closed form.}  The pattern $1, 3, 7, 15, 31, \ldots$ suggests $H_n = 2^n - 1$.

\begin{proof}
\textit{Base step.}  $H_1 = 2^1 - 1 = 1$. \checkmark

\textit{Inductive step.}  Assume $H_k = 2^k - 1$ for some $k \geq 1$.  Then
\[
  H_{k+1} = 2H_k + 1 = 2(2^k - 1) + 1 = 2^{k+1} - 2 + 1 = 2^{k+1} - 1.
\]
\end{proof}

\begin{keyresult}[title={Tower of Hanoi Solution}]
The minimum number of moves to transfer $n$ disks is $H_n = 2^n - 1$.
\end{keyresult}

\begin{remark}[Reve's Puzzle]
If there are $k \geq 4$ pegs instead of 3, the problem is called \emph{Reve's Puzzle}.  The
$k = 4$ case was solved in 2014; the general formula for $k > 4$ remains unknown.
\end{remark}

\subsection{Compound Interest}
\label{subsec:compound-interest}

\begin{example}[Compound Interest]
A deposit of \$10{,}000 earns $11\%$ annual interest compounded annually.  How much is in the
account after $n$ years?
\end{example}

Let $P_n$ be the account balance after $n$ years.  The initial condition is $P_0 = 10{,}000$.

\textbf{Recurrence.}
\[
  P_n = P_{n-1} + 0.11\,P_{n-1} = 1.11\,P_{n-1}.
\]

Iterating this relation:
\[
  P_n = (1.11)^n P_0.
\]

For $n = 30$: $P_{30} = (1.11)^{30} \times 10{,}000 \approx \$228{,}922.97$.

\subsection{Graduate Student Job Problem and Fibonacci Numbers}
\label{subsec:fibonacci}

\begin{example}[Graduate Student Job Problem]
Each professor (except those in their first year) graduates one student per year who immediately
becomes a professor.  First-year professors produce no graduates.  There are no retirements.
Starting with one professor in year 1, how many professors are there in year $n$?
\end{example}

Let $f_n$ be the number of professors in year $n$.  Computing initial terms:
\[
  f_0 = 0,\; f_1 = 1,\; f_2 = 1,\; f_3 = 2,\; f_4 = 3,\; f_5 = 5.
\]

\textbf{Recurrence.}  In year $n$, the total professorial population consists of:
\begin{itemize}
  \item all $f_{n-1}$ professors who were present last year (no retirements), and
  \item the $f_{n-2}$ new graduates produced by professors who have been around for at least
        two years (i.e.\ those present in year $n-2$ are no longer first-year by year $n-1$).
\end{itemize}

Hence
\[
  f_n = f_{n-1} + f_{n-2}, \qquad f_0 = 0,\; f_1 = 1.
\]

This is precisely the recurrence for the \textbf{Fibonacci numbers}.  The closed-form solution
(to be derived in Handout 5) is

\begin{keyresult}[title={Closed Form for Fibonacci Numbers (Binet's Formula)}]
\[
  f_n = \frac{\varphi^n - \psi^n}{\sqrt{5}}, \qquad
  \varphi = \frac{1+\sqrt{5}}{2}, \quad \psi = \frac{1-\sqrt{5}}{2}.
\]
\end{keyresult}

\begin{examtip}[title={Initial Conditions Matter}]
The recurrence $a_n = a_{n-1} + a_{n-2}$ is shared by many sequences.  Always state initial
conditions: the Fibonacci sequence has $f_0 = 0, f_1 = 1$, while the ``no-two-consecutive-zeros''
bit-string count (Section~\ref{subsec:strings}) has $a_1 = 2, a_2 = 3$.
\end{examtip}

\section{Modelling with Recurrence Relations}
\label{sec:recurrence-modelling}

This section develops two further examples that illustrate the method of conditioning on the
last symbol (or digit) of a combinatorial object.

\subsection{Bit Strings with No Two Consecutive Zeros}
\label{subsec:strings}

\begin{example}[Bit Strings]
\label{ex:bitstrings}
How many binary strings of length $n$ contain no two consecutive $0$s?
\end{example}

Let $a_n$ denote this count.  By enumeration: $a_1 = 2$ (strings: \texttt{0}, \texttt{1}) and
$a_2 = 3$ (strings: \texttt{01}, \texttt{10}, \texttt{11}; \texttt{00} is excluded).

\textbf{Recurrence.}  Consider the last bit of a valid string of length $n$:

\begin{itemize}
  \item \textbf{Last bit is $1$:} The first $n-1$ bits can be any valid string of length $n-1$.
        This contributes $a_{n-1}$ strings.
  \item \textbf{Last bit is $0$:} To avoid two consecutive zeros, the $(n-1)$-th bit must be
        $1$.  The first $n-2$ bits can be any valid string of length $n-2$.
        This contributes $a_{n-2}$ strings.
\end{itemize}

Hence
\[
  a_n = a_{n-1} + a_{n-2}, \qquad a_1 = 2,\; a_2 = 3.
\]

The sequence runs $2, 3, 5, 8, 13, \ldots$ — a Fibonacci-type sequence with shifted initial
conditions.  A closed form will be obtained in Handout~5.

\subsection{Valid Decimal Codewords}
\label{subsec:codewords}

\begin{example}[Codewords with an Even Number of Zeros]
\label{ex:codewords}
A decimal string (digits $0$--$9$) is \emph{valid} if it contains an even number of $0$s.
Let $a_n$ be the number of valid $n$-digit strings.  Find a recurrence for $a_n$.
\end{example}

By direct count: $a_1 = 9$, $a_2 = 82$, $a_3 = 756$.

\textbf{Recurrence.}  Condition on the last digit of a valid $n$-digit string:

\begin{itemize}
  \item \textbf{Last digit is non-zero} (9 choices: $1$--$9$): The first $n-1$ digits must form
        a valid codeword (even number of zeros), giving $9\,a_{n-1}$ strings.
  \item \textbf{Last digit is $0$}: The first $n-1$ digits must contain an \emph{odd} number of
        zeros.  Since all $10^{n-1}$ strings of length $n-1$ partition into those with an even
        number of zeros ($a_{n-1}$) and those with an odd number ($10^{n-1} - a_{n-1}$), this
        case contributes $10^{n-1} - a_{n-1}$ strings.
\end{itemize}

Therefore
\[
  a_n = 9\,a_{n-1} + \bigl(10^{n-1} - a_{n-1}\bigr) = 8\,a_{n-1} + 10^{n-1}, \qquad a_1 = 9.
\]

\begin{examtip}[title={Complementary Counting in Recurrences}]
When conditioning on appending a symbol that changes a parity or other property, count the
``bad'' prefix strings as the complement: (\text{total prefixes}) $-$ (\text{valid prefixes}).
\end{examtip}

\begin{warning}[title={Verifying Initial Conditions}]
Always verify that your recurrence is consistent with the initial conditions you state.
For Example~\ref{ex:codewords}: $a_2 = 8 \times 9 + 10^1 = 72 + 10 = 82$. \checkmark
\end{warning}

\section*{Reading Assignment (Handout 4)}
Rosen (7th edition): Chapter~8.1 (Applications of Recurrence Relations).

% ════════════════════════════════════════════════════════════════════════════
%  HANDOUT 5 — LINEAR HOMOGENEOUS RECURRENCE RELATIONS
% ════════════════════════════════════════════════════════════════════════════

\chapter{Linear Homogeneous Recurrence Relations}
\label{chap:linear-homo}

\section{Linear Homogeneous Recurrence Relations}
\label{sec:lhrr}

Recall from Section~\ref{sec:recurrence-intro} that a \emph{recurrence relation} expresses
each term $a_n$ in terms of earlier terms.  In Handout~4 we set up recurrences by modelling
counting problems; here we develop systematic techniques for \emph{solving} them, starting
with the most structured class: linear homogeneous recurrences with constant coefficients.

% ── 5.1 Definition ─────────────────────────────────────────────────────────
\subsection{Definition and Terminology}
\label{subsec:lhrr-def}

\begin{definition}[Linear Homogeneous Recurrence of Degree $k$]
\label{def:lhrr}
A \textbf{linear homogeneous recurrence relation of degree $k$ with constant coefficients}
is a recurrence of the form
\[
  a_n = c_1 a_{n-1} + c_2 a_{n-2} + \cdots + c_k a_{n-k},
\]
where $c_1, c_2, \ldots, c_k \in \mathbb{R}$ and $c_k \neq 0$.
\end{definition}

The three adjectives each rule out something:
\begin{itemize}
  \item \textbf{Linear}: no term $a_{n-i}$ appears raised to a power other than~$1$, and
        there is no product of two sequence terms.
  \item \textbf{Homogeneous}: the right-hand side contains \emph{only} multiples of previous
        terms — no additive constant or other function of $n$.
  \item \textbf{Constant coefficients}: each $c_i$ is a fixed real number, not a function
        of $n$.
\end{itemize}

\begin{example}[Classifying recurrences]
\label{ex:classify}
\hfill
\begin{itemize}
  \item $f_n = f_{n-1} + f_{n-2}$ \quad \checkmark\ linear homogeneous, degree $2$.
  \item $a_n = a_{n-1} + a_{n-2}^2$ \quad \texttimes\ \emph{not linear} ($a_{n-2}^2$).
  \item $H_n = H_{n-1} + 2$ \quad \texttimes\ \emph{not homogeneous} (additive constant).
  \item $X_n = n \cdot X_{n-1}$ \quad \texttimes\ \emph{not constant coefficients} (factor of $n$).
\end{itemize}
\end{example}

\begin{remark}[Uniqueness of solutions]
Given $k$ initial values $a_0, a_1, \ldots, a_{k-1}$, the recurrence
$a_n = c_1 a_{n-1} + \cdots + c_k a_{n-k}$ determines $a_k, a_{k+1}, \ldots$ uniquely.
This follows from the Principle of Strong Mathematical Induction: each new term is computed
from already-known terms.
\end{remark}

% ── 5.2 The Characteristic Equation ────────────────────────────────────────
\subsection{The Characteristic Equation}
\label{subsec:char-eq}

The key idea for solving a linear homogeneous recurrence is to look for solutions of the
form $a_n = r^n$ for some constant $r \neq 0$.  Substituting into
$a_n = c_1 a_{n-1} + c_2 a_{n-2} + \cdots + c_k a_{n-k}$ gives
\[
  r^n = c_1 r^{n-1} + c_2 r^{n-2} + \cdots + c_k r^{n-k}.
\]
Dividing both sides by $r^{n-k}$ (valid since $r \neq 0$) yields the
\emph{characteristic equation}:

\begin{definition}[Characteristic equation]
\label{def:char-eq}
The \textbf{characteristic equation} of $a_n = c_1 a_{n-1} + \cdots + c_k a_{n-k}$ is
\[
  r^k - c_1 r^{k-1} - c_2 r^{k-2} - \cdots - c_{k-1} r - c_k = 0.
\]
Its roots $r_1, \ldots, r_k$ are called the \textbf{characteristic roots}.
\end{definition}

\begin{example}[Finding characteristic equations]
\label{ex:char-eq}
\hfill
\begin{itemize}
  \item $f_n = f_{n-1} + f_{n-2}$: characteristic equation $r^2 - r - 1 = 0$.
  \item $a_n = a_{n-1} + 3a_{n-2} - 2a_{n-3}$: characteristic equation
        $r^3 - r^2 - 3r + 2 = 0$.
\end{itemize}
\end{example}

% ── 5.3 Degree-2: Two Distinct Roots ───────────────────────────────────────
\subsection{Degree-2 Recurrences: Two Distinct Roots}
\label{subsec:deg2-distinct}

We focus first on degree-$2$ recurrences, $a_n = c_1 a_{n-1} + c_2 a_{n-2}$.  If the
characteristic equation $r^2 - c_1 r - c_2 = 0$ has two \emph{distinct} roots $r_1 \neq r_2$,
then both $a_n = r_1^n$ and $a_n = r_2^n$ are solutions.  Moreover, any linear combination
$h_1 r_1^n + h_2 r_2^n$ is also a solution (linearity of the recurrence).  Theorem~1 states
this is in fact the \emph{complete} family of solutions.

\begin{theorem}[Section 8.2, Theorem~1]
\label{thm:deg2-distinct}
Suppose the characteristic equation of $a_n = c_1 a_{n-1} + c_2 a_{n-2}$ has two distinct
roots $r_1$ and $r_2$.  Then the sequence $(a_n)$ is a solution of the recurrence if and
only if
\[
  a_n = h_1 \cdot r_1^n + h_2 \cdot r_2^n
\]
for some constants $h_1, h_2$.
\end{theorem}

\begin{proof}
\textbf{($\Leftarrow$)} Let $a_n = h_1 r_1^n + h_2 r_2^n$.  Since $r_1, r_2$ are roots of
$r^2 - c_1 r - c_2 = 0$, we have $r_i^2 = c_1 r_i + c_2$ for $i = 1,2$.  Therefore:
\begin{align*}
  c_1 a_{n-1} + c_2 a_{n-2}
    &= c_1(h_1 r_1^{n-1} + h_2 r_2^{n-1}) + c_2(h_1 r_1^{n-2} + h_2 r_2^{n-2}) \\
    &= h_1 r_1^{n-2}(c_1 r_1 + c_2) + h_2 r_2^{n-2}(c_1 r_2 + c_2) \\
    &= h_1 r_1^{n-2} \cdot r_1^2 + h_2 r_2^{n-2} \cdot r_2^2 \\
    &= h_1 r_1^n + h_2 r_2^n = a_n.
\end{align*}

\textbf{($\Rightarrow$)} Suppose $(a_n)$ is any solution with $a_0 = C_0$, $a_1 = C_1$.
Set $h_1$ and $h_2$ to satisfy
\[
  \begin{cases} h_1 + h_2 = C_0, \\ h_1 r_1 + h_2 r_2 = C_1. \end{cases}
\]
Since $r_1 \neq r_2$ this linear system has a unique solution:
\[
  h_1 = \frac{C_1 - C_0 r_2}{r_1 - r_2}, \qquad
  h_2 = \frac{C_0 r_1 - C_1}{r_1 - r_2}.
\]
(The denominator $r_1 - r_2 \neq 0$ by distinctness.)  The sequence $b_n = h_1 r_1^n + h_2 r_2^n$
then satisfies the same recurrence \emph{and} the same initial conditions as $(a_n)$.  By
uniqueness of solutions, $a_n = b_n$ for all $n$.
\end{proof}

\begin{keyresult}[title={Solving a Degree-2 Recurrence (Distinct Roots)}]
Given $a_n = c_1 a_{n-1} + c_2 a_{n-2}$ with initial values $a_0, a_1$:
\begin{enumerate}
  \item Write the characteristic equation $r^2 - c_1 r - c_2 = 0$ and find roots $r_1 \neq r_2$.
  \item Write $a_n = h_1 r_1^n + h_2 r_2^n$.
  \item Substitute $n = 0$ and $n = 1$ to form a $2 \times 2$ linear system in $h_1, h_2$.
  \item Solve for $h_1$ and $h_2$.
\end{enumerate}
\end{keyresult}

\begin{warning}[title={Distinctness is Essential}]
Theorem~\ref{thm:deg2-distinct} requires $r_1 \neq r_2$.  When the characteristic equation
has a \emph{repeated} root, the formula $h_1 r^n + h_2 r^n = (h_1+h_2)r^n$ collapses to a
one-parameter family — not enough to match two initial conditions.  Use
Theorem~\ref{thm:deg2-repeated} in that case.
\end{warning}

% ── 5.4 Examples (distinct roots) ──────────────────────────────────────────
\subsection{Examples: Distinct Roots}
\label{subsec:deg2-distinct-examples}

\begin{example}[Section 8.2, Example 3]
\label{ex:sec82-ex3}
Solve $a_n = a_{n-1} + 2a_{n-2}$ with $a_0 = 2$, $a_1 = 7$.

\textbf{Solution.}  The characteristic equation is $r^2 - r - 2 = 0$, which factors as
$(r-2)(r+1) = 0$, giving roots $r_1 = 2$ and $r_2 = -1$.  Thus $a_n = h_1 \cdot 2^n + h_2 \cdot (-1)^n$.
Applying initial conditions:
\[
  \begin{cases} a_0 = 2: & h_1 + h_2 = 2, \\ a_1 = 7: & 2h_1 - h_2 = 7. \end{cases}
  \implies h_1 = 3,\quad h_2 = -1.
\]
Therefore $\boxed{a_n = 3 \cdot 2^n - (-1)^n}$.
\end{example}

\begin{example}[Fibonacci Numbers --- Section 8.2, Example 4]
\label{ex:fibonacci-closed}
Solve $f_n = f_{n-1} + f_{n-2}$ with $f_0 = 0$, $f_1 = 1$.

\textbf{Solution.}  The characteristic equation is $r^2 - r - 1 = 0$, with roots
\[
  r_1 = \frac{1+\sqrt{5}}{2} \quad\text{(the golden ratio } \phi\text{)}, \qquad
  r_2 = \frac{1-\sqrt{5}}{2}.
\]
So $f_n = h_1 \left(\tfrac{1+\sqrt{5}}{2}\right)^n + h_2 \left(\tfrac{1-\sqrt{5}}{2}\right)^n$.
From $f_0 = 0$: $h_1 + h_2 = 0$, so $h_2 = -h_1$.
From $f_1 = 1$: $h_1 \tfrac{1+\sqrt{5}}{2} - h_1 \tfrac{1-\sqrt{5}}{2} = h_1 \sqrt{5} = 1$,
giving $h_1 = 1/\sqrt{5}$.  Therefore
\[
  \boxed{f_n = \frac{1}{\sqrt{5}}\left(\frac{1+\sqrt{5}}{2}\right)^n
              - \frac{1}{\sqrt{5}}\left(\frac{1-\sqrt{5}}{2}\right)^n.}
\]
This is \textbf{Binet's formula}: despite involving irrational numbers, $f_n$ is always a
non-negative integer.
\end{example}

% ── 5.5 Degree-2: Repeated Root ─────────────────────────────────────────────
\subsection{Degree-2 Recurrences: Repeated Root}
\label{subsec:deg2-repeated}

When the characteristic equation $r^2 - c_1 r - c_2 = 0$ has a \emph{repeated} root
$r_0$ (i.e.\ $r^2 - c_1 r - c_2 = (r - r_0)^2$), the second linearly independent solution
is $n \cdot r_0^n$.

\begin{theorem}[Section 8.2, Theorem~2]
\label{thm:deg2-repeated}
Suppose the characteristic equation of $a_n = c_1 a_{n-1} + c_2 a_{n-2}$ has only one
(repeated) root $r_0$.  Then the sequence $(a_n)$ is a solution if and only if
\[
  a_n = h_1 \cdot r_0^n + h_2 \cdot n \cdot r_0^n
\]
for some constants $h_1, h_2$.
\end{theorem}

\begin{examtip}[title={Recognising a Repeated Root}]
If the characteristic equation factors as $(r - r_0)^2 = 0$, use the
$h_1 r_0^n + h_2 n r_0^n$ form.  A quick check: for $r^2 - c_1 r - c_2 = 0$, the discriminant
$c_1^2 + 4c_2 = 0$ signals a repeated root.
\end{examtip}

\begin{example}[Plants reproducing once]
\label{ex:plants}
Each plant produces one new plant during the first year of its life, then never again; all
plants live forever.  Starting with one plant in year~1, how many plants $p_n$ are there in
year $n$?

\textbf{Setting up the recurrence.}  In year $n$ we have all plants from the previous year
($p_{n-1}$) plus the newly born plants, which equals the number of plants that were
\emph{at least one year old} last year, i.e.\ $p_{n-1} - p_{n-2}$.  So
\[
  p_n = p_{n-1} + (p_{n-1} - p_{n-2}) = 2p_{n-1} - p_{n-2},
  \qquad p_0 = 0,\; p_1 = 1.
\]

\textbf{Solving.}  Characteristic equation: $r^2 - 2r + 1 = (r-1)^2 = 0$, so $r_0 = 1$.
General solution: $p_n = h_1 \cdot 1^n + h_2 \cdot n \cdot 1^n = h_1 + h_2 n$.

Applying initial conditions:
\[
  p_0 = 0 \implies h_1 = 0; \qquad p_1 = 1 \implies h_2 = 1.
\]
Therefore $\boxed{p_n = n}$ — precisely as expected, since each year exactly one new plant
enters.
\end{example}

\begin{example}[Section 8.2, Example 5]
\label{ex:sec82-ex5}
Solve $a_n = 6a_{n-1} - 9a_{n-2}$ with $a_0 = 1$, $a_1 = 6$.

\textbf{Solution.}  Characteristic equation: $r^2 - 6r + 9 = (r-3)^2 = 0$, so $r_0 = 3$.
Thus $a_n = h_1 \cdot 3^n + h_2 \cdot n \cdot 3^n$.

From initial conditions:
\[
  \begin{cases} a_0 = 1: & h_1 = 1, \\ a_1 = 6: & 3h_1 + 3h_2 = 6 \implies h_2 = 1. \end{cases}
\]
Therefore $a_n = 3^n + n \cdot 3^n = \boxed{(n+1) \cdot 3^n}$.
\end{example}

% ── 5.6 Higher Degree: Distinct Roots ──────────────────────────────────────
\subsection{Higher-Degree Recurrences: Distinct Roots}
\label{subsec:highdeg-distinct}

The approach extends naturally to degree $k$.

\begin{theorem}[Section 8.2, Theorem~3]
\label{thm:degk-distinct}
Suppose the characteristic equation of
$a_n = c_1 a_{n-1} + c_2 a_{n-2} + \cdots + c_k a_{n-k}$
has $k$ distinct roots $r_1, r_2, \ldots, r_k$.  Then $(a_n)$ is a solution if and only if
\[
  a_n = h_1 r_1^n + h_2 r_2^n + \cdots + h_k r_k^n
\]
for some constants $h_1, \ldots, h_k$.
\end{theorem}

Given $k$ initial conditions $a_0, \ldots, a_{k-1}$, setting $n = 0, 1, \ldots, k-1$
produces a Vandermonde linear system whose unique solution gives $h_1, \ldots, h_k$.

\begin{example}[Section 8.2, Example 6]
\label{ex:sec82-ex6}
Solve $a_n = 6a_{n-1} - 11a_{n-2} + 6a_{n-3}$ with $a_0 = 2$, $a_1 = 5$, $a_2 = 15$.

\textbf{Solution.}  Characteristic equation: $r^3 - 6r^2 + 11r - 6 = 0$.
Testing $r = 1$: $1 - 6 + 11 - 6 = 0$~\checkmark, so $(r-1)$ is a factor.
Dividing: $(r-1)(r^2 - 5r + 6) = (r-1)(r-2)(r-3) = 0$, giving $r_1 = 1$, $r_2 = 2$,
$r_3 = 3$.

Thus $a_n = h_1 \cdot 1^n + h_2 \cdot 2^n + h_3 \cdot 3^n$.  Applying initial conditions:
\[
  \begin{cases} h_1 + h_2 + h_3 = 2, \\ h_1 + 2h_2 + 3h_3 = 5, \\ h_1 + 4h_2 + 9h_3 = 15. \end{cases}
  \implies h_1 = 1,\; h_2 = -1,\; h_3 = 2.
\]
Therefore $\boxed{a_n = 1 - 2^n + 2 \cdot 3^n}$.
\end{example}

% ── 5.7 Repeated Roots (Optional) ──────────────────────────────────────────
\subsection{Repeated Roots in the General Case (Optional)}
\label{subsec:repeated-optional}

When the characteristic equation has roots of higher multiplicity, each root $r_i$ of
multiplicity $m_i$ contributes a block of $m_i$ linearly independent solutions:
$r_i^n,\; n r_i^n,\; n^2 r_i^n,\; \ldots,\; n^{m_i - 1} r_i^n$.

\begin{theorem}[Section 8.2, Theorem~4]
\label{thm:degk-repeated}
Suppose the characteristic equation has $t$ distinct roots $r_1, \ldots, r_t$ with
multiplicities $m_1, \ldots, m_t$ (so $m_1 + \cdots + m_t = k$).  Then $(a_n)$ is a
solution if and only if
\begin{align*}
  a_n = &\;(h_{1,0} + h_{1,1} n + \cdots + h_{1,m_1-1} n^{m_1-1})\,r_1^n \\
        &+ (h_{2,0} + h_{2,1} n + \cdots + h_{2,m_2-1} n^{m_2-1})\,r_2^n \\
        &+ \cdots + (h_{t,0} + h_{t,1} n + \cdots + h_{t,m_t-1} n^{m_t-1})\,r_t^n,
\end{align*}
for constants $h_{i,j}$.
\end{theorem}

\begin{example}[Optional]
\label{ex:triple-root}
Solve $a_n = -3a_{n-1} - 3a_{n-2} - a_{n-3}$ with $a_0 = 1$, $a_1 = -2$, $a_2 = -1$.

\textbf{Solution.}  Characteristic equation: $r^3 + 3r^2 + 3r + 1 = (r+1)^3 = 0$.
One root $r_1 = -1$ with multiplicity $3$.  By Theorem~\ref{thm:degk-repeated}:
\[
  a_n = (h_{1,0} + h_{1,1} n + h_{1,2} n^2)(-1)^n.
\]
Applying initial conditions:
\[
  \begin{cases} h_{1,0} = 1, \\ -(h_{1,0} + h_{1,1} + h_{1,2}) = -2 \implies h_{1,1} + h_{1,2} = 1, \\ h_{1,0} + 2h_{1,1} + 4h_{1,2} = -1 \implies 2h_{1,1} + 4h_{1,2} = -2. \end{cases}
\]
Solving: $h_{1,2} = -2$, $h_{1,1} = 3$.  Therefore
\[
  \boxed{a_n = (1 + 3n - 2n^2)(-1)^n.}
\]
\end{example}

% ── 5.8 Simultaneous Recurrences ───────────────────────────────────────────
\subsection{Simultaneous Recurrence Relations}
\label{subsec:simultaneous}

Sometimes two sequences are coupled: each is defined in terms of the other.  The strategy
is to \emph{eliminate} one sequence to obtain a single higher-degree recurrence, then
back-substitute.

\begin{example}[Simultaneous recurrences]
\label{ex:simultaneous}
Solve
\[
  a_n = a_{n-1} + b_{n-1}, \qquad b_n = a_{n-1} - b_{n-1},
\]
with $a_0 = 1$, $b_0 = 2$.

\textbf{Step 1: eliminate $a$.}  From the second equation, $a_{n-1} = b_n + b_{n-1}$, so
$a_n = b_{n+1} + b_n$.  Substituting into the first equation:
\[
  b_{n+1} + b_n = (b_n + b_{n-1}) + b_{n-1}
  \implies b_{n+1} = 2b_{n-1}
  \implies b_n = 2b_{n-2}.
\]

\textbf{Step 2: solve for $b_n$.}  Characteristic equation: $r^2 - 2 = 0$, roots
$r = \pm\sqrt{2}$.  Thus $b_n = h_1(\sqrt{2})^n + h_2(-\sqrt{2})^n$.

Note $b_1 = a_0 - b_0 = 1 - 2 = -1$.  Applying $b_0 = 2$ and $b_1 = -1$:
\[
  h_1 + h_2 = 2, \quad \sqrt{2}\,h_1 - \sqrt{2}\,h_2 = -1
  \implies h_1 = 1 - \tfrac{1}{2\sqrt{2}},\; h_2 = 1 + \tfrac{1}{2\sqrt{2}}.
\]

\textbf{Step 3: recover $a_n$.}  From $a_n = b_{n+1} + b_n$, substitute the expression
for $b_n$.
\end{example}

\begin{examtip}[title={Solving Simultaneous Recurrences}]
Always begin by using one equation to express one sequence in terms of the other, then
substitute to obtain a \emph{single} recurrence in one unknown.  Once that is solved, recover
the second sequence by back-substitution.
\end{examtip}

% ── 5.9 Connection to Linear Algebra (Optional) ─────────────────────────────
\subsection{Connection to Linear Algebra (Optional)}
\label{subsec:linalg-connection}

There is a beautiful linear-algebraic interpretation of the method above.  Consider again
the degree-$2$ case $a_n = c_1 a_{n-1} + c_2 a_{n-2}$.  Define the vector
$\mathbf{v}_n = \bigl(\begin{smallmatrix} a_n \\ a_{n-1} \end{smallmatrix}\bigr)$.  Then
\[
  \mathbf{v}_n = A\,\mathbf{v}_{n-1}, \quad \text{where } A = \begin{pmatrix} c_1 & c_2 \\ 1 & 0 \end{pmatrix},
\]
so $\mathbf{v}_n = A^{n-1}\mathbf{v}_1$.  If the characteristic equation of the recurrence
has two distinct roots $r_1, r_2$, the matrix $A$ is \emph{diagonalisable}: $A = PDP^{-1}$
where $D = \bigl(\begin{smallmatrix} r_1 & 0 \\ 0 & r_2 \end{smallmatrix}\bigr)$ and
$P$ has the corresponding eigenvectors as columns.  Then $A^{n-1} = PD^{n-1}P^{-1}$, and
computing the first component of $\mathbf{v}_n$ recovers precisely the formula
$a_n = h_1 r_1^n + h_2 r_2^n$.

This perspective explains \emph{why} we get a polynomial-times-exponential form for repeated
roots: a non-diagonalisable (Jordan) block of size $m$ for eigenvalue $r_0$ produces terms
$r_0^n, n r_0^n, \ldots, n^{m-1} r_0^n$.

\section*{Reading Assignment (Handout 5)}
Rosen (7th edition): Chapter~8.2 (Solving Linear Recurrence Relations).

% ============================================================
\chapter{Linear Non-Homogeneous Recurrence Relations}
\label{chap:linear-nonhomo}

\section{Linear Non-Homogeneous Recurrence Relations}
\label{sec:lnhrr}
% ============================================================

In Handout~5 we solved linear \emph{homogeneous} recurrence relations, where every
term depends only on earlier terms and $F(n)=0$.  We now handle the general case where an
additional \emph{driving term} $F(n)\neq 0$ is present.

% ------------------------------------------------------------
\subsection{Definition}
\label{subsec:lnhrr-def}
% ------------------------------------------------------------

\begin{definition}[Linear Non-Homogeneous Recurrence Relation]
\label{def:lnhrr}
A \emph{linear non-homogeneous recurrence relation of degree $k$ with constant
coefficients} has the form
\[
  a_n = c_1 a_{n-1} + c_2 a_{n-2} + \cdots + c_k a_{n-k} + F(n),
\]
where $c_1, \ldots, c_k \in \mathbb{R}$, $c_k \neq 0$, and $F(n) \not\equiv 0$ is a
function of $n$ alone.

The recurrence obtained by setting $F(n) = 0$,
\[
  b_n = c_1 b_{n-1} + c_2 b_{n-2} + \cdots + c_k b_{n-k},
\]
is called the \textbf{associated homogeneous recurrence relation}.
\end{definition}

\begin{keyresult}[title={Key Idea: Two Parts}]
To solve a non-homogeneous recurrence, we need two ingredients:
\begin{enumerate}
  \item A \textbf{particular solution} $a_n^{(p)}$ satisfying the full recurrence.
  \item The \textbf{general solution} $b_n$ of the associated \emph{homogeneous} recurrence (no initial conditions imposed).
\end{enumerate}
Every solution then has the form $a_n^{(p)} + b_n$.
\end{keyresult}

\begin{example}[Identifying associated homogeneous recurrences]
\label{ex:lnhrr-identify}
\hfill
\begin{itemize}
  \item $f_n = f_{n-1} + f_{n-2} + 3$\quad $\Longrightarrow$\quad associated homogeneous: $b_n = b_{n-1} + b_{n-2}$.
  \item $a_n = 3a_{n-1} + 5a_{n-2} + n^2$\quad $\Longrightarrow$\quad associated homogeneous: $b_n = 3b_{n-1} + 5b_{n-2}$.
  \item $a_n = 2a_{n-1} - 5a_{n-3} - 2^n + n^2 - 5$\quad $\Longrightarrow$\quad associated homogeneous: $b_n = 2b_{n-1} - 5b_{n-3}$.
\end{itemize}
Note that $F(n)$ is simply dropped; the coefficients $c_i$ are unchanged.
\end{example}

% ------------------------------------------------------------
\subsection{Structure of the Solution Set}
\label{subsec:lnhrr-structure}
% ------------------------------------------------------------

A key structural observation distinguishes the non-homogeneous from the homogeneous case.

\begin{remark}[Sum of two solutions is \emph{not} a solution]
Suppose $(a_n)$ and $(b_n)$ both satisfy $a_n = c_1 a_{n-1} + \cdots + c_k a_{n-k} + F(n)$.
Then
\begin{align*}
  (a_n + b_n) &= c_1(a_{n-1}+b_{n-1}) + \cdots + c_k(a_{n-k}+b_{n-k}) + 2F(n).
\end{align*}
Since the right-hand side has $2F(n)$, not $F(n)$, the sequence $(a_n + b_n)$ is
\textbf{not} a solution.
\end{remark}

\begin{warning}[title={Warning: Non-Homogeneous Solutions Don't Add}]
  Unlike linear homogeneous recurrences, the sum of two solutions of a non-homogeneous
  recurrence is \emph{not} itself a solution.  Adding a solution of the homogeneous part
  \emph{is} valid (see Theorem~\ref{thm:lnhrr} below).
\end{warning}

However, there is a useful substitute:

\begin{proposition}
\label{prop:part-plus-hom}
If $(a_n)$ is a solution of the non-homogeneous recurrence and $(b_n)$ is a solution of
the \emph{associated homogeneous} recurrence, then $(a_n + b_n)$ is also a solution of the
non-homogeneous recurrence.
\end{proposition}

\begin{proof}
Since $a_n = c_1 a_{n-1} + \cdots + c_k a_{n-k} + F(n)$ and
$b_n = c_1 b_{n-1} + \cdots + c_k b_{n-k}$, adding gives
\[
  (a_n + b_n) = c_1(a_{n-1}+b_{n-1}) + \cdots + c_k(a_{n-k}+b_{n-k}) + F(n). \qedhere
\]
\end{proof}

This leads to the central theorem for this section.

\begin{theorem}[Theorem~5, \S8.2]
\label{thm:lnhrr}
Let $\bigl(a_n^{(p)}\bigr)$ be a \emph{particular} solution of the non-homogeneous
recurrence
\[
  a_n = c_1 a_{n-1} + c_2 a_{n-2} + \cdots + c_k a_{n-k} + F(n).
\]
Then \textbf{every} solution of this recurrence has the form
\[
  a_n = a_n^{(p)} + b_n,
\]
where $(b_n)$ is a solution of the associated homogeneous recurrence
$b_n = c_1 b_{n-1} + \cdots + c_k b_{n-k}$.
(No initial conditions are imposed on $b_n$ at this stage.)
\end{theorem}

\begin{proof}
We already know from Proposition~\ref{prop:part-plus-hom} that $a_n^{(p)} + b_n$ is a
solution.  It remains to show there are no other solutions.

Let $(d_n)$ be any solution of the non-homogeneous recurrence.  Then:
\[
  a_n^{(p)} = c_1 a_{n-1}^{(p)} + \cdots + c_k a_{n-k}^{(p)} + F(n), \qquad
  d_n = c_1 d_{n-1} + \cdots + c_k d_{n-k} + F(n).
\]
Subtracting the first equation from the second:
\[
  d_n - a_n^{(p)} = c_1\!\left(d_{n-1}-a_{n-1}^{(p)}\right) + \cdots +
                    c_k\!\left(d_{n-k}-a_{n-k}^{(p)}\right).
\]
Setting $b_n = d_n - a_n^{(p)}$ shows that $(b_n)$ satisfies the associated homogeneous
recurrence.  Hence $d_n = a_n^{(p)} + b_n$, as claimed.
\end{proof}

\begin{keyresult}[title={Key Result: Theorem 5 (Structure Theorem)}]
Every solution $=$ (one particular solution) $+$ (general solution of associated homogeneous recurrence).

\textbf{Strategy}: (1) Solve the associated homogeneous recurrence to get $b_n$ with free constants.
(2) Find any one particular solution $a_n^{(p)}$.
(3) Write $a_n = a_n^{(p)} + b_n$.
(4) Use initial conditions to pin down the free constants.
\end{keyresult}

% ------------------------------------------------------------
\subsection{Finding a Particular Solution}
\label{subsec:particular}
% ------------------------------------------------------------

The hardest step is guessing $a_n^{(p)}$.  For the two most common forms of $F(n)$,
there are systematic trial forms.

\subsubsection{Polynomial driving term}

\begin{theorem}[Particular solution: polynomial $F$]
\label{thm:particular-poly}
Suppose $F(n) = b_t n^t + b_{t-1}n^{t-1} + \cdots + b_1 n + b_0$ is a polynomial of
degree $t$.
\begin{itemize}
  \item If $1$ is \textbf{not} a root of the characteristic equation of the associated
        homogeneous recurrence, try
        \[a_n^{(p)} = p_t n^t + p_{t-1}n^{t-1} + \cdots + p_1 n + p_0.\]
  \item If $1$ \textbf{is} a root of the characteristic equation, try
        \[a_n^{(p)} = n\cdot(p_t n^t + p_{t-1}n^{t-1} + \cdots + p_1 n + p_0).\]
\end{itemize}
Substitute the trial form into the recurrence and match coefficients to find $p_0, \ldots, p_t$.
\end{theorem}

\begin{example}[Why we multiply by $n$ when $1$ is a root]
\label{ex:root1}
Consider $a_n = a_{n-1} + n$.  The associated homogeneous recurrence is $b_n = b_{n-1}$,
with characteristic equation $r - 1 = 0$, so $r = 1$ is the only root.

Since $F(n) = n$ is a polynomial of degree~1 and $1$ \emph{is} a root, we try
$a_n^{(p)} = n(cn + d) = cn^2 + dn$.  Substituting:
\[
  cn^2 + dn = c(n-1)^2 + d(n-1) + n.
\]
Expanding the right side: $cn^2 - 2cn + c + dn - d + n$.  Matching coefficients:
\[
  n^2{:}\ c = c \quad(\text{automatic}),\qquad
  n^1{:}\ d = -2c + d + 1 \;\Rightarrow\; c = \tfrac{1}{2},\qquad
  n^0{:}\ 0 = c - d \;\Rightarrow\; d = \tfrac{1}{2}.
\]
Hence $a_n^{(p)} = \tfrac{1}{2}n^2 + \tfrac{1}{2}n = \tfrac{n(n+1)}{2}$.

(Had we tried $a_n^{(p)} = cn + d$ without the extra factor of $n$, the coefficient
equations would yield $0 = 1$, a contradiction.)
\end{example}

\subsubsection{Exponential driving term}

\begin{theorem}[Particular solution: exponential $F$]
\label{thm:particular-exp}
Suppose $F(n) = s^n$ for some constant $s \neq 0$.
\begin{itemize}
  \item If $s$ is \textbf{not} a root of the characteristic equation of the associated
        homogeneous recurrence, try $a_n^{(p)} = c \cdot s^n$.
  \item If $s$ \textbf{is} a root of the characteristic equation, try
        $a_n^{(p)} = c \cdot n \cdot s^n$.
\end{itemize}
\end{theorem}

\begin{example}[Why we multiply by $n$ when $s$ is a root]
\label{ex:roots-exp}
Consider $a_n = 5a_{n-1} + 5^n$.  The associated homogeneous recurrence is $b_n = 5b_{n-1}$
with characteristic root $r = 5$.  Since $F(n) = 5^n$ and $s = 5$ \emph{is} a root, we
try $a_n^{(p)} = c \cdot n \cdot 5^n$.  Substituting:
\[
  cn \cdot 5^n = 5 \cdot c(n-1) \cdot 5^{n-1} + 5^n = c(n-1) \cdot 5^n + 5^n.
\]
Dividing both sides by $5^n$: $cn = c(n-1) + 1 = cn - c + 1$, so $c = 1$.

Hence $a_n^{(p)} = n \cdot 5^n$.  (Trying $a_n^{(p)} = c \cdot 5^n$ without the factor
of $n$ leads to $c = c + 1$, a contradiction.)
\end{example}

% ------------------------------------------------------------
\subsection{Worked Examples}
\label{subsec:lnhrr-examples}
% ------------------------------------------------------------

\begin{example}[Example~10, \S8.2 --- Polynomial driving term]
\label{ex:lnhrr-ex10}
Find all solutions of $a_n = 3a_{n-1} + 2n$ with $a_1 = 3$.

\medskip\noindent\textbf{Step 1 --- Solve the associated homogeneous recurrence.}
The associated recurrence is $b_n = 3b_{n-1}$ with characteristic root $r = 3$, giving
\[
  b_n = h \cdot 3^n \quad (h \text{ a free constant}).
\]

\medskip\noindent\textbf{Step 2 --- Find a particular solution.}
Here $F(n) = 2n$ is a polynomial of degree~1.  The characteristic equation is $r = 3$,
so $1$ is not a root.  Try $a_n^{(p)} = cn + d$:
\[
  cn + d = 3\bigl(c(n-1)+d\bigr) + 2n = (3c+2)n + (-3c+3d).
\]
Matching coefficients:
\[
  n^1{:}\ c = 3c + 2 \;\Rightarrow\; c = -1, \qquad
  n^0{:}\ d = -3c + 3d \;\Rightarrow\; 2d = 3c = -3 \;\Rightarrow\; d = -\tfrac{3}{2}.
\]
So $a_n^{(p)} = -n - \tfrac{3}{2}$.

\medskip\noindent\textbf{Step 3 --- General solution (Theorem~\ref{thm:lnhrr}).}
\[
  a_n = -n - \tfrac{3}{2} + h \cdot 3^n.
\]

\medskip\noindent\textbf{Step 4 --- Apply initial condition $a_1 = 3$.}
\[
  3 = -1 - \tfrac{3}{2} + 3h \;\Rightarrow\; 3h = \tfrac{11}{2} \;\Rightarrow\; h = \tfrac{11}{6}.
\]

\medskip\noindent\textbf{Final answer:}
\[
  \boxed{a_n = -n - \tfrac{3}{2} + \tfrac{11}{6}\cdot 3^n.}
\]
\end{example}

\begin{example}[Polynomial driving term, negative coefficient]
\label{ex:lnhrr-neg}
Find all solutions of $a_n = -2a_{n-1} + n + 1$ with $a_0 = 2$.

\medskip\noindent\textbf{Step 1.}
Associated homogeneous: $b_n = -2b_{n-1}$, characteristic root $r = -2$, so
\[
  b_n = h \cdot (-2)^n.
\]

\medskip\noindent\textbf{Step 2.}
$F(n) = n+1$ is degree~1; the characteristic root is $-2 \neq 1$, so try $a_n^{(p)} = cn + d$:
\[
  cn + d = -2\bigl(c(n-1)+d\bigr) + n + 1 = (-2c+1)n + (2c - 2d + 1).
\]
Matching coefficients:
\[
  n^1{:}\ c = -2c+1 \;\Rightarrow\; 3c = 1 \;\Rightarrow\; c = \tfrac{1}{3},\qquad
  n^0{:}\ d = 2c - 2d + 1 \;\Rightarrow\; 3d = \tfrac{2}{3}+1 = \tfrac{5}{3}
         \;\Rightarrow\; d = \tfrac{5}{9}.
\]

\medskip\noindent\textbf{Step 3.}
\[
  a_n = \tfrac{n}{3} + \tfrac{5}{9} + h\cdot(-2)^n.
\]

\medskip\noindent\textbf{Step 4.} Apply $a_0 = 2$:
\[
  2 = 0 + \tfrac{5}{9} + h \;\Rightarrow\; h = 2 - \tfrac{5}{9} = \tfrac{13}{9}.
\]

\medskip\noindent\textbf{Final answer:}
\[
  \boxed{a_n = \frac{n}{3} + \frac{5}{9} + \frac{13}{9}\cdot(-2)^n.}
\]
\end{example}

\begin{example}[Example~11, \S8.2 --- Exponential driving term]
\label{ex:lnhrr-ex11}
Find all solutions of $a_n = 5a_{n-1} - 6a_{n-2} + 7^n$.

\medskip\noindent\textbf{Step 1.}
Associated homogeneous: $b_n = 5b_{n-1} - 6b_{n-2}$.  Characteristic equation:
\[
  r^2 - 5r + 6 = (r-2)(r-3) = 0 \;\Rightarrow\; r_1=2,\; r_2=3.
\]
By Theorem~\ref{thm:deg2-distinct} (Handout~5):
\[
  b_n = h_1 \cdot 2^n + h_2 \cdot 3^n.
\]

\medskip\noindent\textbf{Step 2.}
$F(n) = 7^n$; since $s = 7$ is not a characteristic root, try $a_n^{(p)} = c \cdot 7^n$:
\[
  c \cdot 7^n = 5c \cdot 7^{n-1} - 6c \cdot 7^{n-2} + 7^n.
\]
Dividing by $7^{n-2}$:
\[
  49c = 35c - 6c + 49 \;\Rightarrow\; 20c = 49 \;\Rightarrow\; c = \tfrac{49}{20}.
\]

\medskip\noindent\textbf{Step 3.}
\[
  \boxed{a_n = \frac{49}{20}\cdot 7^n + h_1 \cdot 2^n + h_2 \cdot 3^n,}
\]
where $h_1, h_2$ are determined by two initial conditions.
\end{example}

\begin{examtip}[title={Exam Tip: Choosing the Right Trial Form}]
When $F(n)$ is a polynomial of degree $t$:
\begin{itemize}
  \item Check whether $r = 1$ is a characteristic root of the associated homogeneous recurrence.
  \item If \textbf{no}: try $p_t n^t + \cdots + p_0$ (same degree as $F$).
  \item If \textbf{yes}: multiply by $n$, i.e.\ try $n(p_t n^t + \cdots + p_0)$ (degree $t+1$).
\end{itemize}
When $F(n) = s^n$:
\begin{itemize}
  \item Check whether $s$ is a characteristic root.
  \item If \textbf{no}: try $c \cdot s^n$.
  \item If \textbf{yes}: try $c \cdot n \cdot s^n$.
\end{itemize}
Always substitute your trial form back into the recurrence and match coefficients to find the constants.
\end{examtip}

% ------------------------------------------------------------
\subsection{Recurrence Relations --- Review}
\label{subsec:rr-review}
% ------------------------------------------------------------

Having completed the study of recurrence relations, here is a summary of which theorem to use.

\begin{center}
\renewcommand{\arraystretch}{1.4}
\begin{tabular}{lll}
\hline
\textbf{Theorem} & \textbf{Applies when\ldots} & \textbf{See} \\
\hline
Theorem~1 & Deg-2 homogeneous, two \emph{distinct} characteristic roots & \S\ref{subsec:deg2-distinct} \\
Theorem~2 & Deg-2 homogeneous, one \emph{repeated} characteristic root & \S\ref{subsec:deg2-repeated} \\
Theorem~3 & Deg-$k$ homogeneous, $k$ distinct characteristic roots & \S\ref{subsec:highdeg-distinct} \\
Theorem~5 & Any linear \emph{non-homogeneous} recurrence & \S\ref{sec:lnhrr} \\
\hline
\end{tabular}
\end{center}

\begin{remark}
Theorem~5 always reduces a non-homogeneous problem to (a)~a homogeneous sub-problem
(solved by Theorems~1--3 as appropriate) and (b)~finding one particular solution by the
trial-form method of Section~\ref{subsec:particular}.
\end{remark}

\section*{Reading Assignment (Handout 6)}
Rosen (7th edition): Chapter~8.1 (Applications of Recurrence Relations);
Chapter~8.2 (Solving Linear Recurrence Relations).

\chapter{Graph Theory: Introduction}
\label{chap:graph-theory}

\section{Graph Theory: Introduction}
\label{sec:graph-theory-1}

\subsection{Motivation and Graph Models}
\label{subsec:graph-motivation}

Graphs are mathematical structures that arise naturally whenever we need to model
\emph{relationships} between objects. Real-world examples include:
\begin{itemize}
  \item \textbf{Transit networks} — MRT stations as objects, rail connections as links.
  \item \textbf{Social networks} — users as objects, friendships as links.
  \item \textbf{Molecules} — atoms as objects, chemical bonds as links.
  \item \textbf{Flowcharts / algorithms} — states as objects, transitions as directed links.
\end{itemize}

Each of these structures consists of a \emph{collection of objects} together with \emph{links
between those objects}. The goal of graph theory is to provide a unified framework for
describing and analysing such structures.

A foundational historical problem is the \textbf{Bridges of K\"{o}nigsberg}: in 1735 Euler
asked whether one could take a walk through the city that crosses each of its seven bridges
exactly once. Formalising this question led to the birth of graph theory.

\subsection{Formal Definitions}
\label{subsec:graph-definitions}

\begin{definition}[Graph]
\label{def:graph}
A \textbf{graph} $G = (V, E)$ consists of a nonempty set $V$ of \textbf{vertices} (or
\textbf{nodes}) and a set $E$ of \textbf{edges}. Each edge has either one or two vertices
associated with it, called its \textbf{endpoints}. An edge is said to \textbf{connect} its
endpoints.
\end{definition}

We usually write $n = |V|$ for the number of vertices and $m = |E|$ for the number of edges.

\begin{definition}[Directed Graph]
\label{def:digraph}
A \textbf{directed graph} (or \textbf{digraph}) $(V, E)$ consists of a nonempty vertex set
$V$ and a set of \textbf{directed edges} $E$. Each directed edge $e \in E$ is associated with
an \textbf{ordered pair} $(u, v)$ for some $u, v \in V$, and is said to \emph{start} at $u$
and \emph{end} at $v$.

By contrast, in an \textbf{undirected graph} each edge $e \in E$ is an \textbf{unordered
pair} $\{u, v\}$ for some $u \neq v \in V$. Note $\{u,v\} = \{v,u\}$, whereas
$(u,v) \neq (v,u)$ in general.
\end{definition}

\begin{definition}[Simple Graph and Multigraph]
\label{def:simple-graph}
An undirected graph is a \textbf{simple graph} if
\begin{itemize}
  \item no \emph{self-loops} (edges of the form $\{v, v\}$) are allowed, and
  \item there is at most one edge between any two vertices.
\end{itemize}
For a simple graph on $n$ vertices, $0 \le m \le \binom{n}{2}$.

A \textbf{multigraph} permits self-loops and multiple edges between the same pair of vertices.
\end{definition}

\begin{warning}[title={Convention}]
  Unless stated otherwise, all graphs in this course are \textbf{simple graphs}.
\end{warning}

An important observation: the same graph can be \emph{drawn} in many different ways. The
exact positions of vertices on the page are not part of the mathematical object; only the
vertex set and edge set matter.

\subsection{Graph Terminology}
\label{subsec:graph-terminology}

\begin{definition}[Adjacent, Incident, Endpoints]
\label{def:adjacent}
Let $G = (V, E)$ be a graph. Two vertices $u, v \in V$ are \textbf{adjacent} (or
\textbf{neighbors}) if $\{u, v\} \in E$. If $e = \{u, v\} \in E$, then $e$ is
\textbf{incident} with $u$ and $v$, and both $u$ and $v$ are \textbf{endpoints} of $e$.
\end{definition}

\begin{example}
In the graph with vertices $\{1,2,3,4,5\}$ and edges $\{1,2\},\{1,3\},\{1,5\},\{2,4\},\{3,5\}$:
\begin{itemize}
  \item The adjacent pairs are: $1$--$2$, $1$--$3$, $1$--$5$, $2$--$4$, $3$--$5$.
  \item Edge $e_1 = \{1,2\}$ is incident with vertices $1$ and $2$.
  \item Vertex $1$ is an endpoint of $e_1$, $e_2 = \{1,3\}$, and $e_3 = \{1,5\}$.
\end{itemize}
\end{example}

\begin{definition}[Degree]
\label{def:degree}
The \textbf{degree} of a vertex $v \in V$, written $\deg(v)$, is the number of edges
incident to $v$. The \textbf{degree of the graph} $G$, written $\deg(G)$, is the
\emph{maximum} degree of any vertex in $V$. A vertex of degree $0$ is called
\textbf{isolated}. If all vertices have the same degree $k$, then $G$ is called
\textbf{$k$-regular}.

\textit{Note for multigraphs:} a self-loop at $v$ contributes $2$ to $\deg(v)$.
\end{definition}

\begin{definition}[Subgraph and Induced Subgraph]
\label{def:subgraph}
A \textbf{subgraph} of $G = (V,E)$ is a graph $H = (W, F)$ with $W \subseteq V$ and
$F \subseteq E$. The subgraph $H$ is an \textbf{induced subgraph} if $F$ contains
\emph{all} edges from $E$ whose both endpoints lie in $W$.
\end{definition}

\subsection{The Handshaking Lemma}
\label{subsec:handshaking}

\begin{theorem}[Handshaking Lemma]
\label{thm:handshaking}
Let $G = (V, E)$ be a graph with $m$ edges. Then
\[
  2m = \sum_{v \in V} \deg(v).
\]
\end{theorem}

\begin{proof}
Each edge $\{u,v\} \in E$ contributes exactly $1$ to $\deg(u)$ and $1$ to $\deg(v)$, hence
$2$ in total to the sum $\sum_{v \in V} \deg(v)$. Summing over all $m$ edges gives $2m$.
\end{proof}

\begin{keyresult}[title={Key Result: Handshaking Lemma}]
  In any graph, $2m = \sum_{v \in V} \deg(v)$. Equivalently, the sum of all vertex degrees
  is always even, and equals twice the number of edges.
\end{keyresult}

\begin{theorem}
\label{thm:even-odd-degree}
In every graph $G$, the number of vertices of odd degree is even.
\end{theorem}

\begin{proof}
Partition $V$ into $V_0$ (even-degree vertices) and $V_1$ (odd-degree vertices). By
Theorem~\ref{thm:handshaking},
\[
  \underbrace{2m}_{\text{even}}
  = \underbrace{\sum_{v \in V_0} \deg(v)}_{\text{even (sum of even integers)}}
  + \sum_{v \in V_1} \deg(v).
\]
Hence $\sum_{v \in V_1} \deg(v)$ must also be even. Since each term is odd, the number of
terms $|V_1|$ must be even.
\end{proof}

\begin{example}
How many edges are in a $6$-regular graph on $10$ vertices?

\textit{Solution.} Let $m$ be the number of edges. By the Handshaking Lemma,
$2m = 10 \times 6 = 60$, so $m = 30$.
\end{example}

\begin{examtip}[title={Exam Tip: Applying the Handshaking Lemma}]
  For a $k$-regular graph on $n$ vertices, $2m = nk$, so $m = nk/2$. If $nk$ is odd,
  no such graph exists (enter \textbf{DNE}). Specifically, a $k$-regular graph on $n$ vertices
  exists only when $nk$ is even.
\end{examtip}

\subsection{Special Graphs}
\label{subsec:special-graphs}

\begin{definition}[Complete Graph $K_n$]
\label{def:complete-graph}
The \textbf{complete graph} $K_n$ on $n$ vertices is the unique simple graph in which every
pair of distinct vertices is connected by an edge. It has $\binom{n}{2}$ edges and is
$(n-1)$-regular.
\end{definition}

\begin{definition}[Cycle Graph $C_n$]
\label{def:cycle}
The \textbf{cycle graph} $C_n$ ($n \ge 3$) on vertices $v_1, \ldots, v_n$ has edge set
\[
  \bigl\{\{v_1,v_2\},\{v_2,v_3\},\ldots,\{v_{n-1},v_n\},\{v_n,v_1\}\bigr\}.
\]
Every vertex of $C_n$ has degree $2$, so $C_n$ is $2$-regular.
\end{definition}

\begin{remark}
$C_2$ would require two edges between the same pair of vertices, so it is not a simple graph.
The smallest cycle is $C_3$.
\end{remark}

\subsection{Graph Isomorphism}
\label{subsec:isomorphism}

\begin{definition}[Isomorphic Graphs]
\label{def:isomorphism}
Two undirected graphs $G_1 = (V_1, E_1)$ and $G_2 = (V_2, E_2)$ are \textbf{isomorphic}
if there exists a bijection $f : V_1 \to V_2$ such that for all $u, v \in V_1$,
\[
  \{u,v\} \in E_1 \iff \{f(u), f(v)\} \in E_2.
\]
Intuitively, isomorphic graphs are structurally identical — they differ only in how their
vertices are labelled. Graphs that are not isomorphic are called \textbf{non-isomorphic}.
\end{definition}

Checking isomorphism by brute force requires examining $n!$ bijections for $n$-vertex graphs,
which is infeasible for large $n$. Whether an efficient (polynomial-time) algorithm for graph
isomorphism exists is still an open problem in computer science.

A more practical approach is to use \textbf{graph invariants} — properties preserved by
any isomorphism. If two graphs differ on any invariant, they are non-isomorphic. Common
invariants include:
\begin{itemize}
  \item Number of vertices ($n$) and number of edges ($m$).
  \item The \textbf{degree sequence}: the list of vertex degrees sorted in non-decreasing order.
  \item Number of vertices of each given degree.
\end{itemize}

\begin{warning}[title={Matching Invariants is Necessary, Not Sufficient}]
  Two graphs can have identical degree sequences yet still be non-isomorphic. Matching
  invariants rules out isomorphism when they differ, but does not confirm it when they agree.
\end{warning}

\begin{example}
To show two graphs $G, H$ are non-isomorphic: find a property that holds for $G$ but not $H$.

For instance, suppose vertex $a$ in $G$ has $\deg(a) = 2$ and both its neighbours have
degree $3$. If no degree-$2$ vertex in $H$ is adjacent to two degree-$3$ vertices, then
$G \not\cong H$.
\end{example}

\begin{example}
To show two graphs $G, H$ are isomorphic: exhibit an explicit bijection $f : V(G) \to V(H)$
and verify that $\{u,v\} \in E(G) \iff \{f(u),f(v)\} \in E(H)$ for every pair.
\end{example}

\section*{Reading Assignment (Handout 7)}
Rosen (7th edition): Chapter~10.1 (Graphs and Graph Models);
Chapter~10.2 (Graph Terminology and Special Types of Graphs);
Chapter~10.3 (Graph Isomorphism).

\chapter{Graph Theory: Traversing Graphs}
\label{chap:traversing-graphs}

\section{Directed Graph Degrees}
\label{sec:digraph-degrees}

\subsection{Initial and Terminal Vertices}
\label{subsec:initial-terminal}

Recall from Definition~\ref{def:digraph} that a directed graph $G = (V,E)$ has edges that are
\emph{ordered pairs} of vertices. We now formalise the notion of direction at individual edges.

\begin{definition}[Initial and Terminal Vertex]
\label{def:initial-terminal}
For any edge $e = (u,v)$ in a directed graph $G$, vertex $u$ is called the
\textbf{initial vertex} of $e$, and vertex $v$ is called the \textbf{terminal vertex} of $e$.
\end{definition}

\begin{definition}[In-Degree and Out-Degree]
\label{def:in-out-degree}
In a directed graph $G = (V,E)$:
\begin{itemize}
  \item The \textbf{in-degree} of a vertex $v$, denoted $\deg^{-}(v)$, is the number of edges
    with $v$ as their terminal vertex.
  \item The \textbf{out-degree} of a vertex $v$, denoted $\deg^{+}(v)$, is the number of edges
    with $v$ as their initial vertex.
\end{itemize}
A self-loop at a vertex contributes one to both the in-degree and the out-degree of that vertex.
\end{definition}

\begin{example}[Directed Graph Degrees]
\label{ex:digraph-degrees}
Consider the directed graph with $V = \{v_1, v_2, v_3, v_4, v_5\}$ and
$E = \{(v_1,v_2),(v_1,v_3),(v_1,v_5),(v_4,v_2),(v_3,v_5)\}$.
The degrees are:
\[
\begin{array}{l@{\qquad}l}
  \deg^{-}(v_1) = 0, & \deg^{+}(v_1) = 3, \\
  \deg^{-}(v_2) = 2, & \deg^{+}(v_2) = 0, \\
  \deg^{-}(v_3) = 1, & \deg^{+}(v_3) = 1, \\
  \deg^{-}(v_4) = 0, & \deg^{+}(v_4) = 1, \\
  \deg^{-}(v_5) = 2, & \deg^{+}(v_5) = 0.
\end{array}
\]
\end{example}

\subsection{Directed Handshaking Lemma}
\label{subsec:directed-handshaking}

The undirected handshaking lemma (Theorem~\ref{thm:handshaking}) counted each edge twice
because every edge has two endpoints. In a directed graph each edge has exactly one initial
vertex and one terminal vertex, so we obtain the following analogue.

\begin{theorem}[Directed Handshaking Lemma]
\label{thm:directed-handshaking}
Let $G = (V,E)$ be a directed graph with $m$ edges. Then
\[
  \sum_{v \in V} \deg^{-}(v) \;=\; \sum_{v \in V} \deg^{+}(v) \;=\; m.
\]
\end{theorem}

\begin{proof}
Each edge $(u,v) \in E$ contributes exactly $1$ to $\deg^{+}(u)$ and exactly $1$ to
$\deg^{-}(v)$. Hence summing either quantity over all vertices counts every edge precisely once.
\end{proof}

\begin{examtip}[title={Exam Tip: Directed vs.\ Undirected Handshaking}]
In an \emph{undirected} graph: $\sum \deg(v) = 2m$.
In a \emph{directed} graph: $\sum \deg^{-}(v) = \sum \deg^{+}(v) = m$.
The factor-of-two difference is the most common source of error.
\end{examtip}

\subsection{Graph Type Classification (Recap)}
\label{subsec:graph-types-recap}

For reference, the four standard graph types differ along three axes.
The classification below is the one used in this course (note that the convention on
self-loops in multigraphs is \emph{not} universal across textbooks).

\begin{center}
\begin{tabular}{lccc}
\toprule
\textbf{Graph Type} & \textbf{Self-Loops} & \textbf{Multiple Edges} & \textbf{Directed} \\
\midrule
Simple Graph           & No  & No  & No  \\
Multigraph             & Yes & Yes & No  \\
Directed Simple Graph  & No  & No  & Yes \\
Directed Multigraph    & Yes & Yes & Yes \\
\bottomrule
\end{tabular}
\end{center}

For the remainder of this chapter the discussion applies to both simple graphs and multigraphs
unless stated otherwise.

%% ----------------------------------------------------------------
\section{Paths and Connectivity}
\label{sec:paths-connectivity}

\subsection{Paths, Simple Paths, and Circuits}
\label{subsec:paths}

\begin{definition}[Path]
\label{def:path}
A \textbf{path} in an undirected graph $G$ from $u$ to $v$ is a sequence of vertices
$(u = w_0, w_1, \dots, w_k, w_{k+1} = v)$ such that $\{w_i, w_{i+1}\} \in E$ for all
$0 \le i \le k$.

A \textbf{path} in a directed graph $G$ from $u$ to $v$ is a sequence of vertices
$(u = w_0, w_1, \dots, w_k, w_{k+1} = v)$ such that $(w_i, w_{i+1}) \in E$ for all
$0 \le i \le k$.
\end{definition}

Key terminology:
\begin{itemize}
  \item Paths \emph{can} contain the same edge several times.
  \item The \textbf{length} of a path is the number of edges it contains (at least $1$).
  \item A \textbf{simple path} is a path that contains no edge more than once.
    In a directed graph, a simple path may include both $(u,v)$ and $(v,u)$, each at most once.
  \item If the start vertex equals the end vertex ($u = v$), the path is called a \textbf{circuit}.
\end{itemize}

\begin{warning}[title={Warning: Definition of ``Simple Path''}]
In this course, a \emph{simple path} means no \textbf{edge} is repeated. Some other
textbooks define a simple path as one in which no \textbf{vertex} is repeated ---
a strictly stronger condition. Be sure to use the definition from this course.
\end{warning}

\begin{example}[Paths in a Graph]
\label{ex:paths}
Consider the undirected graph on vertices $\{a,b,c,d,e,f\}$ with edges
$\{a,b\}$, $\{a,d\}$, $\{a,e\}$, $\{b,c\}$, $\{b,e\}$, $\{b,f\}$, $\{c,e\}$, $\{c,f\}$,
$\{d,e\}$, $\{e,f\}$.
\begin{itemize}
  \item $d$--$a$--$b$--$c$--$f$ is a simple path of length $4$.
  \item $d$--$e$--$c$--$b$--$a$--$d$ is a circuit of length $5$.
  \item $a$--$b$, $e$--$f$ is \emph{not} a path because $\{b,e\}$ may not connect them
    consecutively without traversing edge $\{b,e\}$.
  \item $d$--$a$--$b$--$c$--$f$--$b$--$a$--$e$ is a path but \emph{not} simple because the
    edge $\{a,b\}$ is used twice.
\end{itemize}
\end{example}

\subsection{Connectivity}
\label{subsec:connectivity}

\begin{definition}[Connected Graph]
\label{def:connected}
A graph $G = (V,E)$ (simple or multigraph) is \textbf{connected} if for every pair of
distinct vertices $u, v \in V$, there is a path between $u$ and $v$. It is called
\textbf{disconnected} otherwise.

A disconnected graph can be partitioned into several maximal connected subgraphs called
\textbf{connected components}.
\end{definition}

\begin{theorem}
\label{thm:simple-path-connected}
There is a simple path between every pair of distinct vertices of a connected graph.
\end{theorem}

\begin{proof}
Let $u, v$ be distinct vertices in a connected graph $G$. Since $G$ is connected, at least one
path from $u$ to $v$ exists. Let $P$ be a path from $u$ to $v$ of minimum length.

Suppose $P$ is not simple. Then some edge $e = \{x,y\}$ appears at least twice, so either $x$
or $y$ (or both) is visited at least twice. Between the two occurrences of the repeated vertex
we can delete the intermediate sub-path, yielding a shorter path from $u$ to $v$. This
contradicts the minimality of $P$, so $P$ must be simple.
\end{proof}

\subsection{Strong and Weak Connectivity}
\label{subsec:strong-weak}

In a directed graph it is possible that a path exists from $u$ to $v$ but not from $v$ to $u$
(for example, in any directed path graph).

\begin{definition}[Strongly and Weakly Connected]
\label{def:strongly-weakly}
A directed graph $G = (V,E)$ is \textbf{strongly connected} if for every pair of vertices
$u, v \in V$ there is a directed path from $u$ to $v$ \emph{and} a directed path from $v$
to $u$.

A directed graph $G$ is \textbf{weakly connected} if its \textbf{underlying undirected graph}
(obtained by ignoring the directions of all edges) is connected.
\end{definition}

Every strongly connected graph is also weakly connected, but the converse does not hold.

\begin{example}
\label{ex:strong-weak}
A directed cycle on $n$ vertices is strongly connected: from any vertex one can reach any
other by following the cycle. A directed path on $n \ge 2$ vertices is weakly connected but
not strongly connected, since one cannot travel from the terminal vertex back to the initial
vertex.
\end{example}

%% ----------------------------------------------------------------
\section{Euler Circuits and Euler Paths}
\label{sec:euler}

\subsection{The K\"{o}nigsberg Bridges Problem}
\label{subsec:konigsberg}

In 1735 Euler asked whether one could walk through the city of K\"{o}nigsberg crossing each
of its seven bridges exactly once. Modelling the four land masses as vertices and the bridges
as edges produces a multigraph; the question becomes: \emph{does there exist a simple path
that uses every edge?} Euler's negative answer was the first result in graph theory.

\subsection{Definitions}
\label{subsec:euler-def}

\begin{definition}[Euler Circuit and Euler Path]
\label{def:euler}
Let $G$ be a connected graph.
\begin{itemize}
  \item An \textbf{Euler circuit} in $G$ is a simple circuit containing every edge of $G$.
  \item An \textbf{Euler path} in $G$ is a simple path containing every edge of $G$.
\end{itemize}
\end{definition}

Recall that ``simple'' here means no edge is repeated. An Euler circuit is necessarily also a
circuit (starts and ends at the same vertex), while an Euler path need not return to its
starting vertex.

\subsection{Euler Circuit Characterisation}
\label{subsec:euler-circuit}

\begin{theorem}[Euler Circuit Theorem]
\label{thm:euler-circuit}
A connected graph $G$ with at least two vertices has an Euler circuit if and only if every
vertex of $G$ has even degree.
\end{theorem}

\begin{proof}
\textbf{($\Rightarrow$) Only-if direction.}
Suppose $G$ has an Euler circuit $C = (v_1, v_2, \dots, v_m)$ with $v_m = v_1$. For the
starting/ending vertex $v_1$, the first edge $(v_1, v_2)$ and the last edge $(v_{m-1}, v_m)$
each contribute $1$ to $\deg(v_1)$, giving an even contribution of $2$. For any intermediate
occurrence $v_i$ with $1 < i < m$, the edges $(v_{i-1}, v_i)$ and $(v_i, v_{i+1})$ each
contribute $1$ to $\deg(v_i)$. Since $C$ is simple and uses every edge, the total degree of
each vertex is even.

\medskip
\textbf{($\Leftarrow$) If direction.}
Assume every vertex of $G$ has even degree. We construct an Euler circuit as follows.

\emph{Step 1}: Start at an arbitrary vertex $v_1$ and greedily extend a simple path by choosing
unused edges until no unused edge is available at the current vertex. The path must terminate
(the graph has finitely many edges).

\emph{Claim}: The resulting path is a circuit, i.e.\ it returns to $v_1$.

\emph{Proof of claim}: Suppose the path ends at $v_k \ne v_1$. Removing the path's edges from
$G$, every intermediate vertex $v_i$ ($1 < i < k$) loses an even number of incident edges
(one in, one out per visit). But $v_k$, being the endpoint and not the start, loses an
\emph{odd} number. Since $\deg(v_k)$ was even, $v_k$ still has at least one unused edge ---
contradicting the assumption that the algorithm stopped. Hence $v_k = v_1$ and the path is a
circuit $C_1$.

\emph{Step 2}: If $C_1$ uses every edge we are done. Otherwise, remove all edges of $C_1$ from
$G$; every vertex in the remaining graph still has even degree (a circuit removes an even
number of edges at every vertex). Find a vertex $v_i$ on $C_1$ that has an unused edge and
repeat Step~1 from $v_i$ to obtain a new circuit $C_2$ in the remaining graph.

\emph{Step 3}: Splice $C_2$ into $C_1$ at $v_i$ to form a larger circuit. Repeat Steps~1--3
until all edges are included, producing an Euler circuit.
\end{proof}

\begin{keyresult}[title={Key Result: Euler Circuit Test}]
A connected graph has an Euler circuit $\iff$ \textbf{every vertex has even degree}.
This is extremely easy to check and the constructive proof also yields an algorithm for
finding the circuit.
\end{keyresult}

\subsection{Euler Path Characterisation}
\label{subsec:euler-path}

\begin{theorem}[Euler Path Theorem]
\label{thm:euler-path}
A connected graph $G$ has an Euler path but \emph{not} an Euler circuit if and only if $G$
has exactly two vertices of odd degree.
\end{theorem}

\begin{proof}
\textbf{($\Rightarrow$) Only-if direction.}
Suppose $G$ has an Euler path $P = (v_1, \dots, v_m)$ with $v_1 \ne v_m$ (since there is no
Euler circuit). The path leaves $v_1$ one more time than it enters, so $\deg(v_1)$ is odd.
Similarly $\deg(v_m)$ is odd. Every other vertex is entered and left equally often, so it has
even degree. Hence exactly two vertices have odd degree.

\medskip
\textbf{($\Leftarrow$) If direction.}
Suppose $G$ has exactly two vertices of odd degree, say $u$ and $v$. Add a new edge $\{u,v\}$
to obtain $G'$. Now every vertex of $G'$ has even degree, so by Theorem~\ref{thm:euler-circuit}
$G'$ has an Euler circuit. Removing the added edge $\{u,v\}$ from the circuit produces an
Euler path in $G$ from $u$ to $v$.
\end{proof}

\begin{keyresult}[title={Key Result: Euler Path Test}]
A connected graph has an Euler path (but no Euler circuit) $\iff$ it has \textbf{exactly two
vertices of odd degree}. The Euler path must start and end at those two vertices.
\end{keyresult}

\begin{examtip}[title={Exam Tip: Applying the Euler Criteria}]
To determine whether a connected graph has an Euler circuit, an Euler path, or neither:
\begin{enumerate}
  \item Count the number of odd-degree vertices.
  \item $0$ odd-degree vertices $\Rightarrow$ Euler circuit exists.
  \item $2$ odd-degree vertices $\Rightarrow$ Euler path (but no Euler circuit) exists.
  \item Any other count $\Rightarrow$ neither an Euler circuit nor an Euler path exists.
\end{enumerate}
For the K\"{o}nigsberg graph, all four vertices have odd degree ($3, 3, 3, 5$), so neither an
Euler circuit nor an Euler path exists.
\end{examtip}

%% ----------------------------------------------------------------
\section{Hamilton Circuits and Hamilton Paths}
\label{sec:hamilton}

Where Euler paths visit every \emph{edge} exactly once, Hamilton paths visit every
\emph{vertex} exactly once.

\subsection{Definitions}
\label{subsec:hamilton-def}

\begin{definition}[Hamilton Path and Hamilton Circuit]
\label{def:hamilton}
Let $G = (V,E)$ be a graph.
\begin{itemize}
  \item $G$ has a \textbf{Hamilton path} if there is a path that visits every vertex of $G$
    exactly once.
  \item $G$ has a \textbf{Hamilton circuit} if there is a circuit that visits every vertex of $G$
    exactly once (not counting the return to the starting vertex $x_0$).
\end{itemize}
That is, a simple circuit $(x_0, x_1, \dots, x_n, x_0)$ is a Hamilton circuit if the path
$(x_0, x_1, \dots, x_n)$ is a Hamilton path.
\end{definition}

\subsection{Complexity}
\label{subsec:hamilton-complexity}

Unlike the Euler criterion, there is \emph{no} simple necessary-and-sufficient condition
known for the existence of Hamilton paths or circuits. One can test all $n!$ permutations of
the vertex set and check whether consecutive vertices are adjacent, but this brute-force
approach is extremely time-consuming. It is widely believed (though not yet proven) that no
polynomial-time algorithm exists for this problem --- the Hamilton Circuit problem is
\textbf{NP-complete}.

\begin{warning}[title={Warning: Euler $\ne$ Hamilton}]
The two concepts are fundamentally different: Euler deals with \emph{edges}, Hamilton with
\emph{vertices}. An easy degree-based test exists for Euler but not for Hamilton. Do not
conflate the two.
\end{warning}

\section*{Reading Assignment (Handout 8)}
Rosen (7th edition): Chapter~10.4 (Connectivity);
Chapter~10.5 (Euler and Hamilton Paths).

\chapter{Graph Coloring and Planar Graphs}
\label{chap:coloring-planar}

This chapter covers vertex coloring of graphs --- assigning colors to vertices so that no two adjacent vertices share the same color --- as well as bipartite graphs (the ``$2$-colorable'' graphs) and planar graphs, which can be drawn in the plane without edge crossings.

\section{Graph Coloring}
\label{sec:graph-coloring}

\subsection{Vertex Coloring}
\label{subsec:vertex-coloring}

\begin{definition}[Vertex Coloring]
\label{def:k-vertex-coloring}
A \emph{$k$-vertex-coloring} of a graph $G = (V, E)$ is a function
\[
f \colon V \to \{1, \ldots, k\}
\]
such that for all $u, v \in V$ with $(u,v) \in E$, we have $f(u) \neq f(v)$.
\end{definition}

Intuitively, a $k$-vertex-coloring assigns one of $k$ colors to each vertex so that no two adjacent vertices receive the same color.

\subsection{Chromatic Number}
\label{subsec:chromatic-number}

\begin{definition}[Chromatic Number]
\label{def:chromatic-number}
A graph $G$ is \emph{$k$-colorable} if a $k$-vertex-coloring of $G$ exists. The \emph{chromatic number} of $G$, denoted $\chi(G)$, is the smallest $k$ for which $G$ is $k$-colorable.
\end{definition}

\begin{example}[Chromatic Numbers of Standard Graphs]
\label{ex:chromatic-standard}
\leavevmode
\begin{itemize}
  \item $\chi(K_n) = n$: every pair of vertices is adjacent, so all $n$ vertices need distinct colors.
  \item $\chi(C_n) = 2$ for even $n$: alternate two colors around the cycle.
  \item $\chi(C_n) = 3$ for odd $n$: two colors almost suffice, but the last vertex is adjacent to two vertices of different colors, forcing a third.
  \item $\chi(G) = 1$ if and only if $G$ has no edges (all vertices are isolated).
\end{itemize}
\end{example}

\begin{warning}[title={Warning: Computing $\chi(G)$ Is Hard}]
In general, determining the chromatic number of an arbitrary graph is an NP-hard problem. The examples above are special cases with known structure.
\end{warning}

\section{Bipartite Graphs}
\label{sec:bipartite}

\subsection{Definition and Examples}
\label{subsec:bipartite-def}

\begin{definition}[Bipartite Graph]
\label{def:bipartite}
A simple graph $G = (V, E)$ is \emph{bipartite} if its vertex set $V$ can be partitioned into two disjoint sets $V_L$ and $V_R$ such that every edge in $G$ connects a vertex in $V_L$ to a vertex in $V_R$. The pair $(V_L, V_R)$ is called a \emph{bipartition} of~$G$.
\end{definition}

Equivalently, no edge connects two vertices within the same partition. Bipartite graphs arise naturally in modelling relationships between two distinct classes of objects (e.g.\ students and courses, jobs and machines).

\subsection{Bipartiteness and 2-Colorability}
\label{subsec:bipartite-coloring}

\begin{theorem}[Bipartite $\iff$ 2-Colorable]
\label{thm:bipartite-coloring}
A graph $G = (V, E)$ is bipartite if and only if $\chi(G) \leq 2$.
\end{theorem}

\begin{proof}
$(\Leftarrow)$\; If $\chi(G) = 1$, then $G$ has no edges, so any bipartition works. If $\chi(G) = 2$, let $V_L$ be the set of vertices assigned color~1 and $V_R$ those assigned color~2. Then $V_L \cap V_R = \varnothing$ and $V = V_L \cup V_R$. Since vertices within $V_L$ share a color (and likewise within $V_R$), no edge connects two vertices in the same part, so $(V_L, V_R)$ is a valid bipartition.

$(\Rightarrow)$\; Suppose $G$ is bipartite with bipartition $(V_L, V_R)$. Assign color~1 to all vertices in $V_L$ and color~2 to all vertices in $V_R$. Since every edge connects $V_L$ to $V_R$, adjacent vertices always receive different colors. If $E = \varnothing$ then $\chi(G) = 1$; otherwise $\chi(G) = 2$. In either case $\chi(G) \leq 2$.
\end{proof}

\subsection{Odd Cycle Characterisation}
\label{subsec:odd-cycle}

\begin{theorem}[Bipartite Characterisation]
\label{thm:bipartite-odd-cycle}
A graph $G$ is bipartite if and only if $G$ contains no odd cycle as a subgraph.
\end{theorem}

\begin{proof}[Proof sketch]
$(\Rightarrow)$\; If $G$ contains an odd cycle $C_{2k+1}$, then $\chi(C_{2k+1}) = 3$, so $C_{2k+1}$ is not $2$-colorable. Adding vertices and edges to an odd cycle cannot make the resulting graph $2$-colorable, so $G$ is not bipartite.

$(\Leftarrow)$\; Suppose $G$ is not bipartite. Starting from an arbitrary vertex $v$, place all vertices reachable via even-length paths in one set and those via odd-length paths in another. Since the partition fails, some vertex is reachable by both an even-length and an odd-length path from $v$, yielding an odd cycle.
\end{proof}

\begin{keyresult}[title={Key Result: Bipartite Testing}]
To test bipartiteness, run a BFS/DFS and attempt a $2$-coloring. The graph is bipartite if and only if no conflict arises, which is equivalent to having no odd cycle.
\end{keyresult}

\subsection{Complete Bipartite Graphs}
\label{subsec:complete-bipartite}

\begin{definition}[Complete Bipartite Graph]
\label{def:complete-bipartite}
The \emph{complete bipartite graph} $K_{m,n}$ has vertex set $V = V_1 \cup V_2$ with $|V_1| = m$, $|V_2| = n$, and edge set $E = V_1 \times V_2$: every vertex in $V_1$ is adjacent to every vertex in $V_2$, and there are no edges within $V_1$ or $V_2$.
\end{definition}

The graph $K_{m,n}$ has $m + n$ vertices and $mn$ edges. Familiar examples include $K_{3,3}$ (the ``utilities graph'') and $K_{2,5}$.

\section{Planar Graphs}
\label{sec:planar}

\subsection{Definition}
\label{subsec:planar-def}

\begin{definition}[Planar Graph]
\label{def:planar}
A graph is \emph{planar} if it can be drawn in the plane so that no edges cross (except at shared endpoints). Such a crossing-free drawing is called a \emph{planar representation}.
\end{definition}

Note that planarity is a property of the graph itself, not of any particular drawing. A graph that looks like it has crossings may still be planar if an alternative crossing-free drawing exists.

\begin{example}[$K_4$ Is Planar]
\label{ex:k4-planar}
The complete graph $K_4$ appears to have edge crossings in its standard drawing, but it admits a planar representation by moving one vertex outside the triangle formed by the other three.
\end{example}

\begin{example}[$K_{3,3}$ Is Not Planar]
\label{ex:k33-nonplanar}
The complete bipartite graph $K_{3,3}$ cannot be drawn without crossings, no matter how the vertices are arranged. We prove this formally below using Corollary~\ref{cor:planar-no-triangles}.
\end{example}

\subsection{Regions and Euler's Formula}
\label{subsec:euler-formula}

A planar representation divides the plane into contiguous areas called \emph{regions} (or \emph{faces}). The unbounded exterior area is called the \emph{infinite region}. The number of edges bordering a region $R$ is its \emph{degree}, denoted $\deg(R)$; an edge appearing twice on the boundary contributes~$2$.

\begin{proposition}[Sum of Region Degrees]
\label{prop:region-degree-sum}
In any planar representation of a graph, the sum of the degrees of all regions equals twice the number of edges:
\[
\sum_{R} \deg(R) = 2m.
\]
\end{proposition}

This is the ``face'' analogue of the Handshaking Lemma (Theorem~\ref{thm:handshaking}).

\begin{theorem}[Euler's Formula]
\label{thm:euler-formula}
Let $G$ be a connected planar simple graph with $n$ vertices, $m$ edges, and $r$ regions in any planar representation. Then
\[
r = m - n + 2.
\]
\end{theorem}

The proof proceeds by induction on $m$. For the base case $m = 0$, the graph is a single vertex ($n = 1$) with one region (the infinite region), giving $1 = 0 - 1 + 2$. The inductive step considers adding an edge to a connected graph: if the new edge connects to a new vertex, $n$ and $m$ each increase by~$1$ but $r$ stays the same; if it connects two existing vertices, $m$ and $r$ each increase by~$1$ but $n$ stays the same. In both cases the formula is preserved.

\begin{keyresult}[title={Key Result: Euler's Formula}]
For a connected planar simple graph: $r = m - n + 2$. This is the fundamental tool for proving edge bounds and non-planarity results.
\end{keyresult}

\subsection{Edge Bounds for Planar Graphs}
\label{subsec:edge-bounds}

\begin{corollary}[Edge Bound]
\label{cor:planar-edge-bound}
If $G$ is a connected planar simple graph with $n \geq 3$ vertices and $m$ edges, then
\[
m \leq 3n - 6.
\]
\end{corollary}

\begin{proof}
Every region has degree at least $3$ (since the graph is simple with $n \geq 3$), so $2m = \sum_R \deg(R) \geq 3r$, giving $r \leq \frac{2}{3}m$. Substituting into Euler's formula:
\[
2 = n + r - m \leq n + \tfrac{2}{3}m - m = n - \tfrac{1}{3}m,
\]
which rearranges to $m \leq 3n - 6$.
\end{proof}

\begin{corollary}[Edge Bound Without Triangles]
\label{cor:planar-no-triangles}
If $G$ is a connected planar simple graph with $n \geq 3$ vertices, $m$ edges, and no cycles of length~$3$, then
\[
m \leq 2n - 4.
\]
\end{corollary}

\begin{proof}
Without triangles, every region has degree at least $4$, so $2m \geq 4r$, giving $r \leq \frac{m}{2}$. Substituting into Euler's formula:
\[
2 = n + r - m \leq n + \tfrac{1}{2}m - m = n - \tfrac{1}{2}m,
\]
which rearranges to $m \leq 2n - 4$.
\end{proof}

\begin{examtip}[title={Exam Tip: Choosing the Right Corollary}]
Use the $m \leq 3n - 6$ bound for general graphs, but switch to $m \leq 2n - 4$ when the graph has no triangles (e.g.\ bipartite graphs, since all cycles in bipartite graphs have even length $\geq 4$).
\end{examtip}

Both corollaries extend to disconnected graphs: if $G$ has $k$ connected components, add $k-1$ edges to connect them without introducing crossings, then apply the connected-case bound. Since $m \leq m + k - 1 \leq 3n - 6$ (or $2n - 4$), the same inequalities hold.

\subsection{Non-Planarity of $K_5$ and $K_{3,3}$}
\label{subsec:nonplanar-k5-k33}

\begin{example}[$K_5$ Is Not Planar]
\label{ex:k5-nonplanar}
The graph $K_5$ has $n = 5$ vertices and $m = \binom{5}{2} = 10$ edges. By Corollary~\ref{cor:planar-edge-bound}, planarity requires $m \leq 3(5) - 6 = 9$. Since $10 > 9$, $K_5$ is not planar.
\end{example}

\begin{example}[$K_{3,3}$ Is Not Planar]
\label{ex:k33-nonplanar-proof}
The graph $K_{3,3}$ has $n = 6$ vertices and $m = 9$ edges. Being bipartite, $K_{3,3}$ contains no triangles (odd cycles). By Corollary~\ref{cor:planar-no-triangles}, planarity requires $m \leq 2(6) - 4 = 8$. Since $9 > 8$, $K_{3,3}$ is not planar.

Note that Corollary~\ref{cor:planar-edge-bound} alone would give $9 \leq 3(6) - 6 = 12$, which is satisfied --- so the triangle-free bound is essential here.
\end{example}

\begin{warning}[title={Warning: Edge Bounds Are Necessary, Not Sufficient}]
Satisfying $m \leq 3n - 6$ does \emph{not} guarantee planarity. The corollaries provide necessary conditions only: violating them proves non-planarity, but satisfying them does not prove planarity.
\end{warning}

\subsection{Kuratowski's Theorem (Optional)}
\label{subsec:kuratowski}

An \emph{elementary subdivision} of a graph replaces an edge $\{u,v\}$ with a new vertex $w$ and two edges $\{u,w\}$, $\{w,v\}$. Two graphs are \emph{homeomorphic} if they can be obtained from the same graph by sequences of elementary subdivisions.

\begin{theorem}[Kuratowski's Theorem]
\label{thm:kuratowski}
A graph is non-planar if and only if it contains a subgraph homeomorphic to $K_5$ or $K_{3,3}$.
\end{theorem}

Equivalently, $G$ is non-planar if and only if one can find vertices of $G$ corresponding to the vertices of $K_5$ or $K_{3,3}$, connected by internally vertex-disjoint paths of $G$ corresponding to the edges.

\section{Map Coloring and the Four-Color Theorem}
\label{sec:four-color}

The classical map coloring problem asks: given a map divided into contiguous regions, how many colors are needed to color the regions so that no two adjacent regions (sharing a boundary segment, not just a point) receive the same color? This reduces directly to coloring the dual planar graph, where each region becomes a vertex and edges connect adjacent regions.

\begin{theorem}[Four-Color Theorem]
\label{thm:four-color}
Every simple planar graph is $4$-colorable.
\end{theorem}

The problem was posed in 1852. The first accepted proof, by Appel and Haken in 1976, relied on computer-assisted case analysis checking hundreds of unavoidable configurations. Subsequent proofs by Robertson et al.\ (1997) and a formal verification by Gonthier (2005) confirmed the result with smaller case sets and machine-checked logic, respectively.

\begin{keyresult}[title={Key Result: Four-Color Theorem}]
$\chi(G) \leq 4$ for every planar graph $G$. This is tight: some planar graphs (e.g.\ $K_4$) require exactly $4$ colors.
\end{keyresult}

\section*{Reading Assignment (Handout 9)}
Rosen (7th edition): Chapter~10.7 (Planarity); Chapter~10.8 (Coloring).

\chapter{Trees and Spanning Trees}
\label{chap:trees}

\section{Trees}
\label{sec:trees}

\subsection{Definition and Basic Properties}
\label{subsec:tree-def}

\begin{definition}[Tree and Forest]
\label{def:tree}
A \textbf{tree} is a connected graph that does not have a cycle as a subgraph.
A \textbf{forest} is a graph that consists of several trees.
\end{definition}

Forests need not be connected; the connected components of a forest are precisely the individual trees that comprise it.

\begin{definition}[Leaf]
\label{def:leaf}
In a tree $T$, a vertex with degree one is called a \textbf{leaf}.
\end{definition}

Trees and forests enjoy a number of useful structural properties.

\begin{proposition}[Tree Properties]
\label{prop:tree-properties}
Let $T$ be a tree (or forest). Then:
\begin{enumerate}[label=(\roman*)]
  \item $T$ is planar, and its planar representation has exactly one region.
  \item $T$ is bipartite.
  \item If $T$ has $n \geq 2$ vertices, then $T$ has at least two leaves.
  \item Removing a leaf $v$ and the edge incident with $v$ from $T$ still yields a tree.
\end{enumerate}
\end{proposition}

Property~(i) follows immediately from Euler's formula (Theorem~\ref{thm:euler-formula}): a tree with $n$ vertices has $n-1$ edges, so $r = m - n + 2 = (n-1) - n + 2 = 1$. Property~(ii) follows because trees contain no odd cycles (in fact, no cycles at all), which by Theorem~\ref{thm:bipartite-odd-cycle} makes them bipartite.

\subsection{Edges in a Tree}
\label{subsec:tree-edges}

\begin{theorem}[Edge Count]
\label{thm:tree-edges}
A tree with $n$ vertices has exactly $n - 1$ edges.
\end{theorem}

\begin{proof}
We proceed by induction on the number of vertices~$n$.

\textbf{Base case.} When $n = 1$, the tree is a single isolated vertex with $0 = 1 - 1$ edges.

\textbf{Inductive hypothesis.} Assume every tree with $k$ vertices has $k - 1$ edges.

\textbf{Inductive step.} Consider a tree $T$ with $k + 1$ vertices. Since $k + 1 \geq 2$, the tree has at least two leaves; let $v$ be one of them, and let $e$ be the unique edge incident with~$v$. By Proposition~\ref{prop:tree-properties}(iv), $T - \{v\}$ is a tree with $k$ vertices, and by the hypothesis it has $k - 1$ edges. Adding $v$ and $e$ back gives $T$ with $(k - 1) + 1 = k$ edges, completing the induction.
\end{proof}

\begin{keyresult}[title={Alternative Definition of a Tree}]
A connected graph with $n$ vertices and $n - 1$ edges is a tree.
\end{keyresult}

\section{Rooted Trees}
\label{sec:rooted-trees}

\subsection{Definitions and Terminology}
\label{subsec:rooted-tree-def}

\begin{definition}[Rooted Tree]
\label{def:rooted-tree}
A \textbf{rooted tree} is a tree in which one vertex is designated as the \textbf{root} and every edge is directed away from the root.
\end{definition}

In a rooted tree, a vertex that is not the root and has degree one is called a \textbf{leaf}. Vertices that have children are called \textbf{internal vertices}. Note that the root itself may or may not be an internal vertex---if the tree consists of a single vertex, the root has no children and is not internal.

\begin{definition}[Level and Height]
\label{def:level-height}
The \textbf{level} of a vertex $u$ in a rooted tree is the length of the unique path from the root to $u$. The level of the root is defined to be zero. The \textbf{height} of a rooted tree is the maximum level among all its vertices.
\end{definition}

\begin{definition}[Parent, Child, Ancestor, Descendant]
\label{def:parent-child}
Given a rooted tree $T$ and a vertex $v$ other than the root:
\begin{itemize}
  \item The \textbf{parent} of $v$ is the vertex adjacent to $v$ that is one level closer to the root.
  \item The \textbf{children} of $v$ are those vertices adjacent to $v$ that are one level farther from the root.
  \item The \textbf{ancestors} of $v$ are all vertices on the path from the root to $v$, excluding $v$ itself.
  \item The \textbf{descendants} of $v$ are those vertices that have $v$ as an ancestor.
\end{itemize}
The root has no parent, and every non-root vertex has a unique parent.
\end{definition}

\begin{definition}[Subtree]
\label{def:subtree}
If $u$ is a vertex of a rooted tree, the \textbf{subtree rooted at $u$} is the subgraph consisting of $u$, all its descendants, and all edges incident to these descendants.
\end{definition}

\begin{example}[Rooted Tree Terminology]
\label{ex:rooted-tree}
Consider a rooted tree with root $r$, children $\{a, b, c\}$ of $r$, children $\{d, e\}$ of $a$, child $\{f\}$ of $b$, and children $\{g, h\}$ of $c$. Then:
\begin{itemize}
  \item $\mathrm{Parent}(a) = r$, \; $\mathrm{Parent}(d) = a$, \; $r$ has no parent.
  \item $\mathrm{Children}(r) = \{a, b, c\}$, \; $\mathrm{Children}(a) = \{d, e\}$, \; $d$ has no children.
  \item $\mathrm{Ancestor}(c) = \{r\}$, \; $\mathrm{Ancestor}(h) = \{r, c\}$, \; $r$ has no ancestors.
  \item $\mathrm{Descendant}(r) = \{a, b, c, d, e, f, g, h\}$, \; $\mathrm{Descendant}(c) = \{g, h\}$.
\end{itemize}
The leaves are $\{d, e, f, g, h\}$, the internal vertices are $\{r, a, b, c\}$, and the height is~$2$.
\end{example}

\subsection{Binary Trees}
\label{subsec:binary-trees}

\begin{definition}[Binary Tree, Full Binary Tree, Balanced Binary Tree]
\label{def:binary-tree}
\leavevmode
\begin{itemize}
  \item A rooted tree is a \textbf{binary tree} if every vertex has at most two children.
  \item A binary tree is a \textbf{full binary tree} if every internal vertex has exactly two children.
  \item A binary tree of height $h$ is a \textbf{balanced binary tree} if all its leaves are at level $h$ or $h - 1$.
\end{itemize}
\end{definition}

\begin{theorem}[Maximum Leaves in a Binary Tree]
\label{thm:binary-tree-leaves}
There are at most $2^h$ leaves in a binary tree $T$ of height $h$.
\end{theorem}

\begin{proof}
We induct on the height $h$.

\textbf{Base case.} When $h = 1$, the root has at most $2 = 2^1$ children, which are all leaves.

\textbf{Inductive hypothesis.} Assume the claim holds for all binary trees of height at most $k$.

\textbf{Inductive step.} Consider a binary tree $T$ of height $h = k + 1$. The root has at most two children $u$ and $v$; let $T_1$ and $T_2$ be the subtrees rooted at $u$ and $v$, respectively. Each subtree has height at most $k$, so by the hypothesis each contains at most $2^k$ leaves. Since every leaf of $T$ belongs to either $T_1$ or $T_2$, the total number of leaves in $T$ is at most $2^k + 2^k = 2^{k+1}$.
\end{proof}

\subsection{Quick-Fire Properties (True or False)}
\label{subsec:tree-tf}

The following true/false questions help test understanding of tree definitions:
\begin{enumerate}[label=(\alph*)]
  \item A tree with $n$ edges has $n + 1$ vertices. \hfill \textbf{True}
  \item A forest with $n$ vertices must have $n - 1$ edges. \hfill \textbf{False} (a forest with multiple components has fewer)
  \item The chromatic number of a tree is always two or less. \hfill \textbf{True} (trees are bipartite)
  \item An undirected simple graph with $n$ vertices and $n - 1$ edges is always connected. \hfill \textbf{False} (e.g.\ an isolated vertex plus a triangle)
  \item Removing any vertex $v$ from a tree $T$ always leaves a tree. \hfill \textbf{False} (removing an internal vertex disconnects)
  \item A rooted tree with $n$ vertices has $n - 1$ edges. \hfill \textbf{True}
  \item In a rooted tree, the root is always an internal vertex. \hfill \textbf{False} (a single-vertex tree has the root as a leaf)
  \item Leaves of a rooted tree have no children. \hfill \textbf{True}
  \item A rooted tree of height $h$ having $2^h$ vertices is a full binary tree. \hfill \textbf{False}
  \item If $v$ has degree~$2$ in a rooted tree $T$, then $T$ has at least two leaves. \hfill \textbf{False} ($v$ could be the root with one child, the other edge going to its parent---e.g.\ a path on three vertices rooted at the middle vertex has $v$ as root with degree~2 but this is fine; however, a path $r$--$a$--$b$ rooted at $r$ gives $v = a$ with degree~$2$ yet only leaf $b$)
\end{enumerate}

\section{Spanning Trees}
\label{sec:spanning-trees}

\subsection{Definition}
\label{subsec:spanning-tree-def}

\begin{definition}[Spanning Tree]
\label{def:spanning-tree}
Let $G$ be a connected simple graph. A \textbf{spanning tree} of $G$ is a subgraph of $G$ that is a tree containing all vertices of $G$.
\end{definition}

A spanning tree of a graph with $n$ vertices therefore has exactly $n - 1$ edges (by Theorem~\ref{thm:tree-edges}). A connected graph may have many different spanning trees; the task of finding one with desirable properties motivates much of the next handout.

\subsection{Edge-Weighted Graphs}
\label{subsec:edge-weighted}

\begin{definition}[Edge-Weighted Graph]
\label{def:edge-weighted}
An \textbf{edge-weighted graph} is a triple $G = (V, E, W)$, where $V$ is the vertex set, $E$ is the edge set (so $(V, E)$ is a graph), and
$W \colon E \to \mathbb{R}$ is a function assigning a \textbf{weight} to each edge.
\end{definition}

Edge weights can represent many quantities: distance or travel time between vertices, cost of using or maintaining a connection, capacity (e.g.\ pipe flow), failure probability, and so on. Many fundamental graph problems---including the minimum spanning tree and shortest path problems---are formulated on edge-weighted graphs.

\subsection{Minimum Spanning Trees}
\label{subsec:mst}

\begin{definition}[Minimum Spanning Tree]
\label{def:mst}
Given an edge-weighted graph $G$, a \textbf{minimum spanning tree} (MST) of $G$ is a spanning tree whose total edge weight $\sum_{e \in T} W(e)$ is as small as possible among all spanning trees of $G$.
\end{definition}

\begin{example}[Finding an MST by Enumeration]
\label{ex:mst-enumerate}
Consider a graph on six vertices with edges weighted as follows: $W(a,d)=4$, $W(a,c)=2$, $W(d,c)=1$, $W(d,f)=7$, $W(c,e)=8$, $W(f,e)=2$, $W(f,b)=3$, $W(e,b)=6$, $W(c,f)=10$. Among all spanning trees, one can verify by enumeration that the MST uses edges $\{(d,c), (a,c), (f,e), (f,b), (a,d)\}$ (or an equivalent set) with minimum total weight. Algorithmic methods (Kruskal's and Prim's algorithms) for efficiently finding MSTs will be covered in the next handout.
\end{example}

\begin{examtip}[title={Exam Tip: MST vs.\ Spanning Tree}]
A spanning tree must include \emph{all} vertices but uses only $n-1$ edges. The MST additionally minimises total weight. When asked to ``find a spanning tree,'' any acyclic connected subgraph on all vertices suffices; when asked for an MST, you must justify minimality.
\end{examtip}

\section*{Reading Assignment (Handout 10)}

Rosen (7th edition):
\begin{itemize}
  \item Chapter 11.1: Introduction to Trees
  \item Chapter 11.4: Spanning Trees
\end{itemize}

\chapter{Minimum Spanning Trees and Shortest Paths}
\label{chap:mst-shortest}

This chapter covers algorithmic methods for finding minimum spanning trees (building on the definitions in Section~\ref{subsec:mst}) and shortest paths in weighted graphs.

\section{Algorithms for Minimum Spanning Trees}
\label{sec:mst-algorithms}

\subsection{Greedy Approach}
\label{subsec:greedy-mst}

Recall from Definition~\ref{def:mst} that a minimum spanning tree of an edge-weighted graph $G$ is a spanning tree with the smallest total sum of edge weights. The general template for MST algorithms follows a \emph{greedy} strategy: initialise $T = \emptyset$, and repeatedly select an edge $e \in E$ satisfying some local optimality condition, adding $e$ to $T$ until $T$ is a spanning tree.

A greedy algorithm makes the locally optimal choice at each step (``I want the best thing \emph{right now}''), in contrast to non-greedy (strategic) approaches that may sacrifice short-term gains for a globally optimal solution. Remarkably, for the MST problem, greedy algorithms are provably correct.

\subsection{Kruskal's Algorithm}
\label{subsec:kruskal}

\begin{definition}[Kruskal's Algorithm]
\label{def:kruskal}
Process edges in increasing order of cost (breaking ties arbitrarily), and add each edge to $T$ provided it does not create a cycle.
\end{definition}

The formal pseudocode is as follows.

\begin{algorithm}[H]
\caption{Kruskal's Algorithm}
\label{alg:kruskal}
\begin{algorithmic}[1]
\Require Edge-weighted connected graph $G = (V, E)$
\Ensure Minimum spanning tree $T$
\State $T \gets \emptyset$
\While{$T$ is not a spanning tree of $G$}
  \State Choose $e \in E$ of minimum cost
  \State Remove $e$ from $E$
  \If{$T \cup \{e\}$ does not contain a cycle}
    \State $T \gets T \cup \{e\}$
  \EndIf
\EndWhile
\State \Return $T$
\end{algorithmic}
\end{algorithm}

\begin{example}[Kruskal's Algorithm Trace]
\label{ex:kruskal-trace}
Consider a graph on vertices $\{a, b, c, d, e, f, g\}$ with edge weights: $w(a,d)=5$, $w(c,e)=5$, $w(d,f)=6$, $w(a,b)=7$, $w(b,e)=7$, $w(b,c)=8$, $w(e,f)=8$, $w(d,b)=9$, $w(e,g)=9$, $w(d,e)=15$, $w(f,g)=11$.

Sorting edges by weight: $\{a,d\}(5)$, $\{c,e\}(5)$, $\{d,f\}(6)$, $\{a,b\}(7)$, $\{b,e\}(7)$, $\{b,c\}(8)$, $\{e,f\}(8)$, $\{d,b\}(9)$, $\{e,g\}(9)$, $\{d,e\}(15)$, $\{f,g\}(11)$.

Kruskal's algorithm adds edges in the following order:
\begin{enumerate}
  \item $\{a,d\}$ (weight 5) --- no cycle formed.
  \item $\{c,e\}$ (weight 5) --- no cycle formed.
  \item $\{d,f\}$ (weight 6) --- no cycle formed.
  \item $\{a,b\}$ (weight 7) --- no cycle formed.
  \item $\{b,e\}$ (weight 7) --- no cycle: connects the $\{a,b,d,f\}$ and $\{c,e\}$ components.
  \item $\{b,c\}$ (weight 8) --- \emph{skipped}: would create cycle $b$--$e$--$c$--$b$.
  \item $\{e,g\}$ (weight 9) --- no cycle; all 7 vertices now connected.
\end{enumerate}
The MST has total weight $5+5+6+7+7+9 = 39$ and uses 6 edges ($= |V|-1$).
\end{example}

\subsection{Prim's Algorithm}
\label{subsec:prim}

\begin{definition}[Prim's Algorithm]
\label{def:prim}
Starting from a single vertex, the set $T$ maintained by the algorithm is always a tree. In each iteration, pick the edge with the least cost that connects a vertex in $T$ to a vertex not yet in $T$.
\end{definition}

\begin{algorithm}[H]
\caption{Prim's Algorithm}
\label{alg:prim}
\begin{algorithmic}[1]
\Require Edge-weighted connected graph $G = (V, E)$
\Ensure Minimum spanning tree $T$
\State $T \gets \emptyset$
\State $S \gets \{v_0\}$ for some starting vertex $v_0$
\While{$T$ is not a spanning tree of $G$}
  \State Choose $e = (u,v) \in E$ of minimum cost with $u \in S$, $v \in V \setminus S$
  \State $T \gets T \cup \{e\}$
  \State $S \gets S \cup \{v\}$
\EndWhile
\State \Return $T$
\end{algorithmic}
\end{algorithm}

\begin{example}[Prim's Algorithm Trace]
\label{ex:prim-trace}
Using the same graph as Example~\ref{ex:kruskal-trace}, starting from vertex $a$:
\begin{enumerate}
  \item $S = \{a\}$: cheapest edge leaving $S$ is $\{a,d\}$ (weight 5). Add it; $S = \{a,d\}$.
  \item $S = \{a,d\}$: cheapest crossing edge is $\{d,f\}$ (weight 6). Add it; $S = \{a,d,f\}$.
  \item $S = \{a,d,f\}$: cheapest crossing edge is $\{a,b\}$ (weight 7). Add it; $S = \{a,d,f,b\}$.
  \item $S = \{a,d,f,b\}$: cheapest crossing edge is $\{b,e\}$ (weight 7). Add it; $S = \{a,d,f,b,e\}$.
  \item $S = \{a,d,f,b,e\}$: cheapest crossing edge is $\{c,e\}$ (weight 5). Add it; $S = \{a,d,f,b,e,c\}$.
  \item $S = \{a,d,f,b,e,c\}$: cheapest crossing edge is $\{e,g\}$ (weight 9). Add it; $S = V$.
\end{enumerate}
The MST edges are $\{a,d\}, \{d,f\}, \{a,b\}, \{b,e\}, \{c,e\}, \{e,g\}$ with total weight 39, the same MST found by Kruskal's algorithm (the edge insertion order differs).
\end{example}

\begin{remark}
\label{rem:mst-uniqueness}
When all edge weights are distinct, the MST is unique. When some weights coincide, there may be multiple MSTs, but all have the same total weight. Different tie-breaking choices in Kruskal's or Prim's algorithms may yield different MSTs in this case.
\end{remark}

\section{Correctness of MST Algorithms}
\label{sec:mst-correctness}

Both Kruskal's and Prim's algorithms are provably correct. Their correctness rests on the \emph{cut property} of minimum spanning trees. For simplicity, we assume all edge costs are distinct (so the MST is unique).

\subsection{Cuts and Safe Edges}
\label{subsec:cuts-safe}

\begin{definition}[Cut]
\label{def:cut}
Given a graph $G = (V, E)$, a \emph{cut} is a partition of the vertex set into two disjoint non-empty subsets $(S, V \setminus S)$. The \emph{edges of the cut} are those with one endpoint in $S$ and the other in $V \setminus S$; such edges are said to \emph{cross} the cut.
\end{definition}

\begin{definition}[Safe Edge]
\label{def:safe-edge}
An edge $e$ is a \emph{safe edge} if there exists a cut $(S, V \setminus S)$ such that $e$ is the unique minimum-cost edge crossing that cut.
\end{definition}

\begin{example}[Identifying Safe Edges]
\label{ex:safe-edge}
Consider a graph on vertices $\{a,b,c,d,e,f\}$ with edge weights $w(a,b)=4$, $w(a,c)=6$, $w(b,c)=12$, $w(b,d)=3$, $w(b,e)=7$, $w(c,d)=10$, $w(c,f)=13$, $w(d,e)=5$, $w(d,f)=8$, $w(e,f)=10$, $w(a,d)=15$.

The edge $\{a,c\}$ (weight 6) is safe because for the cut $S = \{a,b,d\}$, $V \setminus S = \{c,e,f\}$, the edges crossing are $\{a,c\}(6)$, $\{b,e\}(7)$, $\{b,c\}(12)$, $\{c,d\}(10)$, $\{d,e\}(5)$... In fact, one must check: edge $\{a,c\}$ is the cheapest edge crossing the cut $S=\{a,b,d,e,f\}$, $V\setminus S=\{c\}$, since the edges crossing this cut are $\{a,c\}(6)$, $\{b,c\}(12)$, $\{c,d\}(10)$, $\{c,f\}(13)$, and $\{a,c\}$ is the cheapest.
\end{example}

\subsection{The Cut Property}
\label{subsec:cut-property}

\begin{theorem}[Cut Property]
\label{thm:cut-property}
Given an undirected connected graph $G$ with distinct edge costs, the set of safe edges in $G$ forms the unique MST of $G$.
\end{theorem}

The proof proceeds via two lemmas.

\begin{lemma}
\label{lem:safe-in-mst}
If $e$ is a safe edge, then every MST contains $e$.
\end{lemma}

\begin{proof}
Suppose for contradiction that $e = \{u,v\}$ is safe but is not in some MST $T$. Since $e$ is safe, there exists a cut $(S, V \setminus S)$ such that $e$ is the unique minimum-cost edge crossing it. Adding $e$ to $T$ creates a cycle $C$ in $T \cup \{e\}$. Since $u \in S$ and $v \in V \setminus S$, traversing $C$ from $u$ to $v$ without using $e$ must cross the cut at least once via some other edge $e' \in T$. Since $e$ is the unique minimum-cost crossing edge, $w(e) < w(e')$. Then $T' = (T \setminus \{e'\}) \cup \{e\}$ is a spanning tree with strictly lower cost than $T$, contradicting the minimality of $T$.
\end{proof}

\begin{lemma}
\label{lem:safe-cover}
The set of safe edges forms a connected graph that covers every vertex.
\end{lemma}

\begin{proof}
Suppose for contradiction that the graph induced by safe edges does not cover all vertices. Let $S$ be a connected component in this graph. Then $V \setminus S \neq \emptyset$, and the minimum-cost edge $e$ crossing the cut $(S, V \setminus S)$ is safe by definition. But then $e$ is a safe edge not in the graph induced by safe edges, a contradiction.
\end{proof}

\begin{keyresult}[title={Key Result: Why Kruskal's and Prim's Work}]
Both algorithms only add \emph{safe} edges. Kruskal's algorithm: when edge $e = \{u,v\}$ is added, letting $S$ be the connected component containing $u$, the edge $e$ is the minimum-cost edge crossing $(S, V \setminus S)$. Prim's algorithm: at each step, the chosen edge is explicitly the minimum-cost edge crossing $(S, V \setminus S)$ where $S$ is the current tree.
\end{keyresult}

\subsection{Correctness of Kruskal's Algorithm}
\label{subsec:kruskal-correct}

\begin{proof}
It suffices to show that every edge added by Kruskal's algorithm is safe. When $e = \{u,v\}$ is added, let $S$ be the connected component containing $u$ in the current forest. Since $e$ does not form a cycle, $v \notin S$. Among all edges crossing the cut $(S, V \setminus S)$, any edge with lower cost than $e$ would have been processed earlier and either added (connecting $S$ to another component, contradicting $v \notin S$) or rejected (for forming a cycle within $S$). Hence $e$ has minimum cost among edges crossing $(S, V \setminus S)$, so $e$ is safe.
\end{proof}

\subsection{Correctness of Prim's Algorithm}
\label{subsec:prim-correct}

\begin{proof}
At each step, let $S$ be the set of vertices connected by the edges in $T$. The algorithm explicitly selects $e$ as the minimum-cost edge crossing the cut $(S, V \setminus S)$, so $e$ is safe. Since $S$ starts as a single vertex and grows by one vertex per iteration, eventually $S = V$ and $T$ is a spanning tree consisting entirely of safe edges.
\end{proof}

\section{Shortest Path Problems}
\label{sec:shortest-paths}

\subsection{Definitions}
\label{subsec:sp-defs}

\begin{definition}[Shortest Path and Distance]
\label{def:shortest-path}
For an edge-weighted graph $G = (V, E, W)$ with non-negative edge weights:
\begin{itemize}
  \item A \emph{shortest path} between vertices $u$ and $v$ is a simple path from $u$ to $v$ whose total edge weight is minimal.
  \item The \emph{length} of a path is the sum of its edge weights.
  \item The \emph{distance} $d(u,v)$ between $u$ and $v$ is the length of a shortest path between them. If no path exists, $d(u,v) = \infty$.
\end{itemize}
\end{definition}

Three variants of the shortest path problem arise:
\begin{enumerate}
  \item \textbf{Single-pair}: Given vertices $s$ and $t$, find a shortest path from $s$ to $t$.
  \item \textbf{Single-source}: Given vertex $s$, find shortest paths from $s$ to all other vertices.
  \item \textbf{All-pairs}: Find shortest paths for every pair of vertices.
\end{enumerate}
We focus on problems (1) and (2), using directed graphs (undirected problems reduce to directed ones by replacing each undirected edge with two directed edges).

\subsection{Algorithm Ideas}
\label{subsec:sp-ideas}

The key insight is to explore vertices in increasing order of distance from the source $s$. Let $d(v)$ denote the shortest path length from $s$ to $v$.

\begin{lemma}[Optimal Substructure]
\label{lem:optimal-substructure}
If $s = v_0 \to v_1 \to \cdots \to v_k = v$ is a shortest path from $s$ to $v$, then for any $1 \le i < k$, the sub-path $s = v_0 \to v_1 \to \cdots \to v_i$ is a shortest path from $s$ to $v_i$.
\end{lemma}

\begin{proof}
Suppose for contradiction that there exists a shorter path $P'$ from $s$ to $v_i$. Then replacing the sub-path $s \to \cdots \to v_i$ with $P'$ and appending $v_i \to \cdots \to v_k$ yields a path from $s$ to $v_k$ with strictly smaller total weight, contradicting the optimality of the original path.
\end{proof}

At the beginning of the $i$-th iteration, the set $S$ stores the $i-1$ closest vertices to $s$.

\begin{proposition}
\label{prop:next-closest}
Let $v$ be the $i$-th closest vertex to $s$, and let $P$ be a shortest path from $s$ to $v$. Then all intermediate vertices of $P$ belong to $S$.
\end{proposition}

\begin{proof}
If $v'$ is an intermediate vertex on $P$, then $d(v') < d(v)$ (since edge weights are non-negative), so $v'$ is among the $i-1$ closest vertices and hence $v' \in S$.
\end{proof}

This means the next closest vertex is always adjacent to $S$. Define, for each $u \in V \setminus S$:
\[
  \pi(u) = \min_{a \in S} \{ d(a) + \ell(a, u) \}
\]
where $\ell(a,u)$ is the weight of edge $(a,u)$ (or $\infty$ if no such edge exists).

\begin{lemma}
\label{lem:pi-equals-d}
If $u$ is the $i$-th closest vertex to $s$, then $\pi(u) = d(u)$.
\end{lemma}

The $i$-th closest vertex to $s$ is therefore $u \in V \setminus S$ minimising $\pi(u)$.

\subsection{Dijkstra's Algorithm}
\label{subsec:dijkstra}

Rather than recomputing $\pi(v)$ from scratch each iteration, we maintain $d(v)$ as a running estimate that is updated incrementally: when vertex $u$ is added to $S$, for each neighbour $v$ of $u$ we set $d(v) \gets \min\{d(v),\; d(u) + \ell(u,v)\}$.

\begin{algorithm}[H]
\caption{Dijkstra's Algorithm}
\label{alg:dijkstra}
\begin{algorithmic}[1]
\Require Directed weighted graph $G = (V, E)$ with non-negative edge weights, source $s$
\Ensure Shortest distances $d(v)$ for all $v \in V$
\State $d(v) \gets \infty$ for all $v \in V$; \quad $d(s) \gets 0$
\State $S \gets \emptyset$
\While{$S \neq V$}
  \State Find $u \in V \setminus S$ with $d(u) = \min_{v \in V \setminus S} d(v)$
  \State $S \gets S \cup \{u\}$
  \For{each $v \in \mathrm{Adj}(u)$}
    \State $d(v) \gets \min\{d(v),\; d(u) + \ell(u,v)\}$
  \EndFor
\EndWhile
\end{algorithmic}
\end{algorithm}

\begin{warning}[title={Warning: Non-Negative Weights Required}]
Dijkstra's algorithm requires all edge weights to be non-negative. If negative weights are present, the greedy selection of the minimum-distance vertex may be incorrect, and algorithms such as Bellman--Ford must be used instead.
\end{warning}

\begin{example}[Dijkstra's Algorithm Trace]
\label{ex:dijkstra-trace}
Consider a graph on vertices $\{a, b, c, d, e, z\}$ with edges: $\ell(a,b)=4$, $\ell(a,c)=2$, $\ell(c,b)=1$, $\ell(b,d)=5$, $\ell(c,d)=8$, $\ell(c,e)=10$, $\ell(d,e)=2$, $\ell(d,z)=6$, $\ell(e,z)=3$. Source: $s = a$.

\begin{center}
\begin{tabular}{c|cccccc|l}
\hline
Round & $d(a)$ & $d(b)$ & $d(c)$ & $d(d)$ & $d(e)$ & $d(z)$ & $S$ \\
\hline
0 & \textbf{0} & $\infty$ & $\infty$ & $\infty$ & $\infty$ & $\infty$ & $\emptyset$ \\
1 & 0 & 4 & \textbf{2} & $\infty$ & $\infty$ & $\infty$ & $\{a\}$ \\
2 & 0 & \textbf{3} & 2 & 10 & 12 & $\infty$ & $\{a, c\}$ \\
3 & 0 & 3 & 2 & \textbf{8} & 12 & $\infty$ & $\{a, c, b\}$ \\
4 & 0 & 3 & 2 & 8 & \textbf{10} & 14 & $\{a, c, b, d\}$ \\
5 & 0 & 3 & 2 & 8 & 10 & \textbf{13} & $\{a, c, b, d, e\}$ \\
6 & 0 & 3 & 2 & 8 & 10 & 13 & $\{a, c, b, d, e, z\}$ \\
\hline
\end{tabular}
\end{center}

The shortest path from $a$ to $z$ has length 13, via $a \to c \to b \to d \to e \to z$.
\end{example}

\begin{examtip}[title={Exam Tip: Dijkstra Table Method}]
When running Dijkstra's algorithm by hand, set up a table with one column per vertex and one row per round. In each round: (1) identify the unvisited vertex with smallest $d$-value (bold it), (2) add it to $S$, (3) update $d$-values for its neighbours. The bold entries trace the order vertices are finalised.
\end{examtip}

\begin{example}[Dijkstra on an Undirected Graph]
\label{ex:dijkstra-undirected}
Consider an undirected graph on $\{a,b,c,d,e,f\}$ with weights:
$w(a,b)=6$, $w(a,c)=3$, $w(b,c)=2$, $w(b,d)=6$, $w(c,d)=10$,
$w(c,e)=4$, $w(d,e)=2$, $w(d,z)=3$, $w(e,z)=6$. Source: $s = a$.

\begin{center}
\begin{tabular}{c|cccccc|l}
\hline
Round & $d(a)$ & $d(b)$ & $d(c)$ & $d(d)$ & $d(e)$ & $d(z)$ & $S$ \\
\hline
0 & \textbf{0} & $\infty$ & $\infty$ & $\infty$ & $\infty$ & $\infty$ & $\emptyset$ \\
1 & 0 & 6 & \textbf{3} & $\infty$ & $\infty$ & $\infty$ & $\{a\}$ \\
2 & 0 & \textbf{5} & 3 & 13 & 7 & $\infty$ & $\{a,c\}$ \\
3 & 0 & 5 & 3 & 11 & \textbf{7} & $\infty$ & $\{a,c,b\}$ \\
4 & 0 & 5 & 3 & \textbf{9} & 7 & 13 & $\{a,c,b,e\}$ \\
5 & 0 & 5 & 3 & 9 & 7 & \textbf{12} & $\{a,c,b,e,d\}$ \\
6 & 0 & 5 & 3 & 9 & 7 & 12 & $\{a,c,b,e,d,z\}$ \\
\hline
\end{tabular}
\end{center}
\end{example}

\section*{Reading Assignment (Handout 11)}

Rosen (7th edition):
\begin{itemize}
  \item Chapter 11.5: Minimum Spanning Trees
  \item Chapter 10.6: Shortest-Path Problems
\end{itemize}

\end{document}
